% ============================================================
% Projekt układu sumatora pełnego (Full Adder)
% z minimalizacją bramek logicznych
% ------------------------------------------------------------
% Dokument główny — kompilacja w Overleaf (pdfLaTeX).
% Skleja pliki 01–06 w całą pracę.
% ============================================================

\documentclass[12pt,a4paper]{article}

\usepackage[utf8]{inputenc}
\usepackage[T1]{fontenc}
\input{glyphtounicode}
\pdfgentounicode=1
\usepackage{cmap}
\usepackage[polish]{babel}
\IfFileExists{lmodern.sty}{\usepackage{lmodern}}{}
\usepackage[margin=2.5cm]{geometry}
\usepackage{amsmath, amssymb}
\usepackage{array, booktabs}
\usepackage{tikz}
\usepackage{circuitikz}
\usepackage{setspace}
\usepackage{hyperref}

\hypersetup{
    pdftitle={Projekt układu sumatora pełnego (Full Adder) z minimalizacją bramek logicznych},
    pdfauthor={[IMIĘ NAZWISKO NUMER INDEKSU]},
    pdfsubject={Elektronika -- projekt zaliczeniowy},
    pdfcreator={LaTeX}
}

\onehalfspacing
\setlength{\parindent}{0em}
\setlength{\parskip}{0.6em}

\begin{document}

% ---------- Strona tytułowa ----------
\begin{titlepage}
    \centering
    {\Large \textbf{ELEKTRONIKA}\par}
    \vspace{1em}
    {\large Prowadzący: prof. dr inż. Zbigniew Frąckiewicz\par}
    \vspace{4em}
    {\LARGE \textbf{Projekt układu sumatora pełnego (Full Adder)\\[0.5em] z minimalizacją bramek logicznych}\par}
    \vspace{4em}
    {\large 18. 09. 2026\par}
    \vfill
    {\large \textbf{Wyższa Szkoła Informatyki i Zarządzania}\\[0.3em]
    Bielsko-Biała, 2025/2026\par}
    \vspace{2em}
    \begin{tabular}{|c|}
        \hline
        [IMIĘ NAZWISKO NUMER INDEKSU] \\
        \hline
    \end{tabular}
\end{titlepage}

% ---------- Spis treści ----------
\tableofcontents
\newpage

% ---------- Treść pracy ----------
% ============================================================
% Punkt 1: Wstęp teoretyczny
% Projekt układu sumatora pełnego (Full Adder)
% z minimalizacją bramek logicznych
% ------------------------------------------------------------
% Plik przeznaczony do włączenia w dokumencie głównym Overleaf:
%   % ============================================================
% Punkt 1: Wstęp teoretyczny
% Projekt układu sumatora pełnego (Full Adder)
% z minimalizacją bramek logicznych
% ------------------------------------------------------------
% Plik przeznaczony do włączenia w dokumencie głównym Overleaf:
%   % ============================================================
% Punkt 1: Wstęp teoretyczny
% Projekt układu sumatora pełnego (Full Adder)
% z minimalizacją bramek logicznych
% ------------------------------------------------------------
% Plik przeznaczony do włączenia w dokumencie głównym Overleaf:
%   \input{01_wstep_teoretyczny}
% Wymagane pakiety w dokumencie głównym:
%   \usepackage[utf8]{inputenc}  \usepackage[T1]{polski}
%   \usepackage{amsmath, amssymb}
% ============================================================

\section{Wstęp teoretyczny}

\subsection{Definicja sumatora cyfrowego}

Sumator cyfrowy jest układem kombinacyjnym służącym do wykonywania operacji arytmetycznego dodawania liczb zapisanych w postaci binarnej. Jego zadaniem jest wyznaczenie sumy wartości podanych na wejścia układu wraz z ewentualnym przeniesieniem powstającym podczas dodawania. W zależności od zastosowania sumator może operować na pojedynczych bitach lub na wielobitowych słowach danych.

W niniejszym projekcie analizowany jest sumator pełny (Full Adder), którego zadaniem jest dodanie trzech bitów: dwóch bitów argumentów oraz bitu przeniesienia z niższego rzędu. Układ posiada trzy wejścia binarne oraz dwa wyjścia informujące o wyniku dodawania.

Dla przyjętego rozwiązania:
\begin{equation*}
A = a_1a_0, \qquad B = b_1b_0,
\end{equation*}
gdzie $a_i$ są bitami pierwszej liczby, natomiast $b_i$ bitami drugiej liczby. W przypadku pojedynczej pozycji sumator pełny realizuje zależność:
\begin{equation*}
S = A + B + C_{in}.
\end{equation*}

Wyjście układu oznaczone jako $S$ (suma) przyjmuje wartość:
\begin{equation*}
S = 1
\end{equation*}
gdy:
\begin{equation*}
a \oplus b \oplus c_{in} = 1,
\end{equation*}
oraz:
\begin{equation*}
S = 0
\end{equation*}
gdy:
\begin{equation*}
a \oplus b \oplus c_{in} = 0.
\end{equation*}

Drugim wyjściem układu jest przeniesienie $C_{out}$, które przyjmuje wartość logiczną 1, gdy co najmniej dwa z trzech sygnałów wejściowych znajdują się w stanie wysokim.

Sumator pełny może być zrealizowany przy wykorzystaniu podstawowych bramek logicznych, takich jak AND, OR oraz NOT. Możliwe jest również wykorzystanie bramek XOR. W projekcie szczególną uwagę poświęcono minimalizacji liczby wykorzystanych bramek logicznych, dzięki czemu możliwe jest uzyskanie prostszego i bardziej efektywnego układu.

\subsection{Zasada działania sumatora pełnego}

Sumator pełny wykonuje dodawanie trzech bitów zgodnie z regułami arytmetyki binarnej. Podczas dodawania pojedynczych bitów może powstać dwucyfrowy wynik, którego starszy bit stanowi przeniesienie do następnej pozycji.

Dla dwóch bitów $a$ i $b$ możliwe są następujące przypadki:
\begin{align*}
0 + 0 &= 00, \\
0 + 1 &= 01, \\
1 + 0 &= 01, \\
1 + 1 &= 10.
\end{align*}

Sumator pełny uwzględnia dodatkowo trzeci składnik, którym jest przeniesienie $c_{in}$ otrzymane z niższego rzędu. Dodawanie trzech bitów może więc wytworzyć wynik dwucyfrowy o postaci $c_{out}s$, gdzie $s$ jest bitem sumy, a $c_{out}$ przeniesieniem do rzędu wyższego.

Działanie układu można opisać za pomocą następujących zależności. Suma bitów wyraża się wzorem:
\begin{equation*}
s = a \oplus b \oplus c_{in},
\end{equation*}
natomiast przeniesienie wynosi:
\begin{equation*}
c_{out} = a \cdot b \vee c_{in} \cdot (a \oplus b).
\end{equation*}

Wprowadzając oznaczenie pomocnicze:
\begin{equation*}
c_1 = a \oplus b,
\end{equation*}
działanie układu można przedstawić jako dwa etapy:
\begin{equation*}
c_1 = a \oplus b, \qquad s = c_1 \oplus c_{in}, \qquad c_{out} = a \cdot b \vee c_1 \cdot c_{in}.
\end{equation*}

Dzięki temu wyjście $s$ będzie równe 1 wyłącznie wtedy, gdy nieparzysta liczba sygnałów wejściowych przyjmuje wartość logiczną 1, natomiast przeniesienie $c_{out}$ pojawi się w sytuacji, gdy wartość sumy argumentów przekroczy pojemność pojedynczego bitu.

Symbol blokowy sumatora pełnego przedstawia Rysunek~\ref{fig:symbol-fa}.

\begin{figure}[h]
    \centering
    \begin{tikzpicture}
        \draw[thick] (0,0) rectangle (3,2.4);
        \node at (1.5,1.9) {\textbf{Full Adder}};
        \draw[thick] (-1,1.8) -- (0,1.8) node[midway, above] {};
        \node[left] at (-1,1.8) {$a$};
        \draw[thick] (-1,1.2) -- (0,1.2);
        \node[left] at (-1,1.2) {$b$};
        \draw[thick] (-1,0.6) -- (0,0.6);
        \node[left] at (-1,0.6) {$c_{in}$};
        \draw[thick] (3,1.5) -- (4,1.5);
        \node[right] at (4,1.5) {$s$};
        \draw[thick] (3,0.9) -- (4,0.9);
        \node[right] at (4,0.9) {$c_{out}$};
    \end{tikzpicture}
    \caption{Symbol blokowy sumatora pełnego (Full Adder)}
    \label{fig:symbol-fa}
\end{figure}

\subsection{Zastosowanie sumatorów cyfrowych}

Sumatory cyfrowe są powszechnie stosowane w układach cyfrowych, w których konieczne jest wykonywanie operacji arytmetycznych na wartościach binarnych.

Przykładowe zastosowania sumatorów obejmują:
\begin{itemize}
    \item procesory i jednostki arytmetyczno-logiczne (ALU),
    \item układy dodawania liczb wielobitowych,
    \item układy mnożenia i dzielenia binarnego,
    \item liczniki i układy inkrementacji wartości,
    \item układy adresowania pamięci,
    \item cyfrowe przetwarzanie sygnałów (DSP),
    \item układy korekcji kodu BCD,
    \item systemy mikroprocesorowe i mikrokontrolerowe.
\end{itemize}

Sumator pełny może być wykorzystany między innymi do realizacji jednej pozycji dodawania liczb wielobitowych. Przykładowo, układ może stanowić pojedyncze ogniwo kaskady sumującej liczby o dowolnej długości słowa.

W większych systemach sumatory pełne są łączone ze sobą w celu realizacji dodawania liczb o większej liczbie bitów, przy czym przeniesienie z każdego rzędu podawane jest na wejście rzędu wyższego.

\subsection{Rodzaje sumatorów}

Sumatory cyfrowe można podzielić ze względu na liczbę obsługiwanych argumentów oraz sposób przekazywania przeniesienia. Najczęściej wyróżnia się trzy podstawowe typy:
\begin{itemize}
    \item Półsumator (Half Adder) -- dodaje dwa bity $a$ i $b$, wyznaczając sumę $s = a \oplus b$ oraz przeniesienie $c = a \cdot b$; nie posiada wejścia przeniesienia z niższego rzędu.
    \item Sumator pełny (Full Adder) -- dodaje trzy bity $a$, $b$ oraz $c_{in}$, wyznaczając sumę oraz przeniesienie do rzędu wyższego; stanowi podstawowy element układów sumujących wielobitowych.
    \item Sumator kaskadowy (Ripple-Carry Adder) -- powstaje przez połączenie kaskadowe sumatorów pełnych, w którym przeniesienie z każdego rzędu podawane jest na wejście przeniesienia rzędu następnego.
\end{itemize}

W praktycznych układach wyróżnia się również bardziej zaawansowane struktury, takie jak sumatory z przyspieszonym przeniesieniem (Carry-Lookahead Adder), w których przeniesienia wyznaczane są równolegle dla wszystkich pozycji, co pozwala zredukować czas propagacji sygnału kosztem zwiększenia złożoności układu.

Ze względu na liczbę bitów obsługiwanych przez pojedynczy stopień można wyróżnić sumatory jedno- oraz wielobitowe. Zwiększenie liczby bitów powoduje wzrost złożoności układu, dlatego istotnym elementem projektowania układów cyfrowych jest ich odpowiednia minimalizacja.

W niniejszym projekcie zastosowany zostanie pojedynczy sumator pełny. Układ będzie dodawał trzy bity wejściowe $a$, $b$ oraz $c_{in}$, a jego wyjścia $s$ oraz $c_{out}$ będą informowały o wartości sumy oraz przeniesienia. Następnie układ zostanie poddany procesowi minimalizacji w celu ograniczenia liczby wymaganych bramek logicznych.

% Wymagane pakiety w dokumencie głównym:
%   \usepackage[utf8]{inputenc}  \usepackage[T1]{polski}
%   \usepackage{amsmath, amssymb}
% ============================================================

\section{Wstęp teoretyczny}

\subsection{Definicja sumatora cyfrowego}

Sumator cyfrowy jest układem kombinacyjnym służącym do wykonywania operacji arytmetycznego dodawania liczb zapisanych w postaci binarnej. Jego zadaniem jest wyznaczenie sumy wartości podanych na wejścia układu wraz z ewentualnym przeniesieniem powstającym podczas dodawania. W zależności od zastosowania sumator może operować na pojedynczych bitach lub na wielobitowych słowach danych.

W niniejszym projekcie analizowany jest sumator pełny (Full Adder), którego zadaniem jest dodanie trzech bitów: dwóch bitów argumentów oraz bitu przeniesienia z niższego rzędu. Układ posiada trzy wejścia binarne oraz dwa wyjścia informujące o wyniku dodawania.

Dla przyjętego rozwiązania:
\begin{equation*}
A = a_1a_0, \qquad B = b_1b_0,
\end{equation*}
gdzie $a_i$ są bitami pierwszej liczby, natomiast $b_i$ bitami drugiej liczby. W przypadku pojedynczej pozycji sumator pełny realizuje zależność:
\begin{equation*}
S = A + B + C_{in}.
\end{equation*}

Wyjście układu oznaczone jako $S$ (suma) przyjmuje wartość:
\begin{equation*}
S = 1
\end{equation*}
gdy:
\begin{equation*}
a \oplus b \oplus c_{in} = 1,
\end{equation*}
oraz:
\begin{equation*}
S = 0
\end{equation*}
gdy:
\begin{equation*}
a \oplus b \oplus c_{in} = 0.
\end{equation*}

Drugim wyjściem układu jest przeniesienie $C_{out}$, które przyjmuje wartość logiczną 1, gdy co najmniej dwa z trzech sygnałów wejściowych znajdują się w stanie wysokim.

Sumator pełny może być zrealizowany przy wykorzystaniu podstawowych bramek logicznych, takich jak AND, OR oraz NOT. Możliwe jest również wykorzystanie bramek XOR. W projekcie szczególną uwagę poświęcono minimalizacji liczby wykorzystanych bramek logicznych, dzięki czemu możliwe jest uzyskanie prostszego i bardziej efektywnego układu.

\subsection{Zasada działania sumatora pełnego}

Sumator pełny wykonuje dodawanie trzech bitów zgodnie z regułami arytmetyki binarnej. Podczas dodawania pojedynczych bitów może powstać dwucyfrowy wynik, którego starszy bit stanowi przeniesienie do następnej pozycji.

Dla dwóch bitów $a$ i $b$ możliwe są następujące przypadki:
\begin{align*}
0 + 0 &= 00, \\
0 + 1 &= 01, \\
1 + 0 &= 01, \\
1 + 1 &= 10.
\end{align*}

Sumator pełny uwzględnia dodatkowo trzeci składnik, którym jest przeniesienie $c_{in}$ otrzymane z niższego rzędu. Dodawanie trzech bitów może więc wytworzyć wynik dwucyfrowy o postaci $c_{out}s$, gdzie $s$ jest bitem sumy, a $c_{out}$ przeniesieniem do rzędu wyższego.

Działanie układu można opisać za pomocą następujących zależności. Suma bitów wyraża się wzorem:
\begin{equation*}
s = a \oplus b \oplus c_{in},
\end{equation*}
natomiast przeniesienie wynosi:
\begin{equation*}
c_{out} = a \cdot b \vee c_{in} \cdot (a \oplus b).
\end{equation*}

Wprowadzając oznaczenie pomocnicze:
\begin{equation*}
c_1 = a \oplus b,
\end{equation*}
działanie układu można przedstawić jako dwa etapy:
\begin{equation*}
c_1 = a \oplus b, \qquad s = c_1 \oplus c_{in}, \qquad c_{out} = a \cdot b \vee c_1 \cdot c_{in}.
\end{equation*}

Dzięki temu wyjście $s$ będzie równe 1 wyłącznie wtedy, gdy nieparzysta liczba sygnałów wejściowych przyjmuje wartość logiczną 1, natomiast przeniesienie $c_{out}$ pojawi się w sytuacji, gdy wartość sumy argumentów przekroczy pojemność pojedynczego bitu.

Symbol blokowy sumatora pełnego przedstawia Rysunek~\ref{fig:symbol-fa}.

\begin{figure}[h]
    \centering
    \begin{tikzpicture}
        \draw[thick] (0,0) rectangle (3,2.4);
        \node at (1.5,1.9) {\textbf{Full Adder}};
        \draw[thick] (-1,1.8) -- (0,1.8) node[midway, above] {};
        \node[left] at (-1,1.8) {$a$};
        \draw[thick] (-1,1.2) -- (0,1.2);
        \node[left] at (-1,1.2) {$b$};
        \draw[thick] (-1,0.6) -- (0,0.6);
        \node[left] at (-1,0.6) {$c_{in}$};
        \draw[thick] (3,1.5) -- (4,1.5);
        \node[right] at (4,1.5) {$s$};
        \draw[thick] (3,0.9) -- (4,0.9);
        \node[right] at (4,0.9) {$c_{out}$};
    \end{tikzpicture}
    \caption{Symbol blokowy sumatora pełnego (Full Adder)}
    \label{fig:symbol-fa}
\end{figure}

\subsection{Zastosowanie sumatorów cyfrowych}

Sumatory cyfrowe są powszechnie stosowane w układach cyfrowych, w których konieczne jest wykonywanie operacji arytmetycznych na wartościach binarnych.

Przykładowe zastosowania sumatorów obejmują:
\begin{itemize}
    \item procesory i jednostki arytmetyczno-logiczne (ALU),
    \item układy dodawania liczb wielobitowych,
    \item układy mnożenia i dzielenia binarnego,
    \item liczniki i układy inkrementacji wartości,
    \item układy adresowania pamięci,
    \item cyfrowe przetwarzanie sygnałów (DSP),
    \item układy korekcji kodu BCD,
    \item systemy mikroprocesorowe i mikrokontrolerowe.
\end{itemize}

Sumator pełny może być wykorzystany między innymi do realizacji jednej pozycji dodawania liczb wielobitowych. Przykładowo, układ może stanowić pojedyncze ogniwo kaskady sumującej liczby o dowolnej długości słowa.

W większych systemach sumatory pełne są łączone ze sobą w celu realizacji dodawania liczb o większej liczbie bitów, przy czym przeniesienie z każdego rzędu podawane jest na wejście rzędu wyższego.

\subsection{Rodzaje sumatorów}

Sumatory cyfrowe można podzielić ze względu na liczbę obsługiwanych argumentów oraz sposób przekazywania przeniesienia. Najczęściej wyróżnia się trzy podstawowe typy:
\begin{itemize}
    \item Półsumator (Half Adder) -- dodaje dwa bity $a$ i $b$, wyznaczając sumę $s = a \oplus b$ oraz przeniesienie $c = a \cdot b$; nie posiada wejścia przeniesienia z niższego rzędu.
    \item Sumator pełny (Full Adder) -- dodaje trzy bity $a$, $b$ oraz $c_{in}$, wyznaczając sumę oraz przeniesienie do rzędu wyższego; stanowi podstawowy element układów sumujących wielobitowych.
    \item Sumator kaskadowy (Ripple-Carry Adder) -- powstaje przez połączenie kaskadowe sumatorów pełnych, w którym przeniesienie z każdego rzędu podawane jest na wejście przeniesienia rzędu następnego.
\end{itemize}

W praktycznych układach wyróżnia się również bardziej zaawansowane struktury, takie jak sumatory z przyspieszonym przeniesieniem (Carry-Lookahead Adder), w których przeniesienia wyznaczane są równolegle dla wszystkich pozycji, co pozwala zredukować czas propagacji sygnału kosztem zwiększenia złożoności układu.

Ze względu na liczbę bitów obsługiwanych przez pojedynczy stopień można wyróżnić sumatory jedno- oraz wielobitowe. Zwiększenie liczby bitów powoduje wzrost złożoności układu, dlatego istotnym elementem projektowania układów cyfrowych jest ich odpowiednia minimalizacja.

W niniejszym projekcie zastosowany zostanie pojedynczy sumator pełny. Układ będzie dodawał trzy bity wejściowe $a$, $b$ oraz $c_{in}$, a jego wyjścia $s$ oraz $c_{out}$ będą informowały o wartości sumy oraz przeniesienia. Następnie układ zostanie poddany procesowi minimalizacji w celu ograniczenia liczby wymaganych bramek logicznych.

% Wymagane pakiety w dokumencie głównym:
%   \usepackage[utf8]{inputenc}  \usepackage[T1]{polski}
%   \usepackage{amsmath, amssymb}
% ============================================================

\section{Wstęp teoretyczny}

\subsection{Definicja sumatora cyfrowego}

Sumator cyfrowy jest układem kombinacyjnym służącym do wykonywania operacji arytmetycznego dodawania liczb zapisanych w postaci binarnej. Jego zadaniem jest wyznaczenie sumy wartości podanych na wejścia układu wraz z ewentualnym przeniesieniem powstającym podczas dodawania. W zależności od zastosowania sumator może operować na pojedynczych bitach lub na wielobitowych słowach danych.

W niniejszym projekcie analizowany jest sumator pełny (Full Adder), którego zadaniem jest dodanie trzech bitów: dwóch bitów argumentów oraz bitu przeniesienia z niższego rzędu. Układ posiada trzy wejścia binarne oraz dwa wyjścia informujące o wyniku dodawania.

Dla przyjętego rozwiązania:
\begin{equation*}
A = a_1a_0, \qquad B = b_1b_0,
\end{equation*}
gdzie $a_i$ są bitami pierwszej liczby, natomiast $b_i$ bitami drugiej liczby. W przypadku pojedynczej pozycji sumator pełny realizuje zależność:
\begin{equation*}
S = A + B + C_{in}.
\end{equation*}

Wyjście układu oznaczone jako $S$ (suma) przyjmuje wartość:
\begin{equation*}
S = 1
\end{equation*}
gdy:
\begin{equation*}
a \oplus b \oplus c_{in} = 1,
\end{equation*}
oraz:
\begin{equation*}
S = 0
\end{equation*}
gdy:
\begin{equation*}
a \oplus b \oplus c_{in} = 0.
\end{equation*}

Drugim wyjściem układu jest przeniesienie $C_{out}$, które przyjmuje wartość logiczną 1, gdy co najmniej dwa z trzech sygnałów wejściowych znajdują się w stanie wysokim.

Sumator pełny może być zrealizowany przy wykorzystaniu podstawowych bramek logicznych, takich jak AND, OR oraz NOT. Możliwe jest również wykorzystanie bramek XOR. W projekcie szczególną uwagę poświęcono minimalizacji liczby wykorzystanych bramek logicznych, dzięki czemu możliwe jest uzyskanie prostszego i bardziej efektywnego układu.

\subsection{Zasada działania sumatora pełnego}

Sumator pełny wykonuje dodawanie trzech bitów zgodnie z regułami arytmetyki binarnej. Podczas dodawania pojedynczych bitów może powstać dwucyfrowy wynik, którego starszy bit stanowi przeniesienie do następnej pozycji.

Dla dwóch bitów $a$ i $b$ możliwe są następujące przypadki:
\begin{align*}
0 + 0 &= 00, \\
0 + 1 &= 01, \\
1 + 0 &= 01, \\
1 + 1 &= 10.
\end{align*}

Sumator pełny uwzględnia dodatkowo trzeci składnik, którym jest przeniesienie $c_{in}$ otrzymane z niższego rzędu. Dodawanie trzech bitów może więc wytworzyć wynik dwucyfrowy o postaci $c_{out}s$, gdzie $s$ jest bitem sumy, a $c_{out}$ przeniesieniem do rzędu wyższego.

Działanie układu można opisać za pomocą następujących zależności. Suma bitów wyraża się wzorem:
\begin{equation*}
s = a \oplus b \oplus c_{in},
\end{equation*}
natomiast przeniesienie wynosi:
\begin{equation*}
c_{out} = a \cdot b \vee c_{in} \cdot (a \oplus b).
\end{equation*}

Wprowadzając oznaczenie pomocnicze:
\begin{equation*}
c_1 = a \oplus b,
\end{equation*}
działanie układu można przedstawić jako dwa etapy:
\begin{equation*}
c_1 = a \oplus b, \qquad s = c_1 \oplus c_{in}, \qquad c_{out} = a \cdot b \vee c_1 \cdot c_{in}.
\end{equation*}

Dzięki temu wyjście $s$ będzie równe 1 wyłącznie wtedy, gdy nieparzysta liczba sygnałów wejściowych przyjmuje wartość logiczną 1, natomiast przeniesienie $c_{out}$ pojawi się w sytuacji, gdy wartość sumy argumentów przekroczy pojemność pojedynczego bitu.

Symbol blokowy sumatora pełnego przedstawia Rysunek~\ref{fig:symbol-fa}.

\begin{figure}[h]
    \centering
    \begin{tikzpicture}
        \draw[thick] (0,0) rectangle (3,2.4);
        \node at (1.5,1.9) {\textbf{Full Adder}};
        \draw[thick] (-1,1.8) -- (0,1.8) node[midway, above] {};
        \node[left] at (-1,1.8) {$a$};
        \draw[thick] (-1,1.2) -- (0,1.2);
        \node[left] at (-1,1.2) {$b$};
        \draw[thick] (-1,0.6) -- (0,0.6);
        \node[left] at (-1,0.6) {$c_{in}$};
        \draw[thick] (3,1.5) -- (4,1.5);
        \node[right] at (4,1.5) {$s$};
        \draw[thick] (3,0.9) -- (4,0.9);
        \node[right] at (4,0.9) {$c_{out}$};
    \end{tikzpicture}
    \caption{Symbol blokowy sumatora pełnego (Full Adder)}
    \label{fig:symbol-fa}
\end{figure}

\subsection{Zastosowanie sumatorów cyfrowych}

Sumatory cyfrowe są powszechnie stosowane w układach cyfrowych, w których konieczne jest wykonywanie operacji arytmetycznych na wartościach binarnych.

Przykładowe zastosowania sumatorów obejmują:
\begin{itemize}
    \item procesory i jednostki arytmetyczno-logiczne (ALU),
    \item układy dodawania liczb wielobitowych,
    \item układy mnożenia i dzielenia binarnego,
    \item liczniki i układy inkrementacji wartości,
    \item układy adresowania pamięci,
    \item cyfrowe przetwarzanie sygnałów (DSP),
    \item układy korekcji kodu BCD,
    \item systemy mikroprocesorowe i mikrokontrolerowe.
\end{itemize}

Sumator pełny może być wykorzystany między innymi do realizacji jednej pozycji dodawania liczb wielobitowych. Przykładowo, układ może stanowić pojedyncze ogniwo kaskady sumującej liczby o dowolnej długości słowa.

W większych systemach sumatory pełne są łączone ze sobą w celu realizacji dodawania liczb o większej liczbie bitów, przy czym przeniesienie z każdego rzędu podawane jest na wejście rzędu wyższego.

\subsection{Rodzaje sumatorów}

Sumatory cyfrowe można podzielić ze względu na liczbę obsługiwanych argumentów oraz sposób przekazywania przeniesienia. Najczęściej wyróżnia się trzy podstawowe typy:
\begin{itemize}
    \item Półsumator (Half Adder) -- dodaje dwa bity $a$ i $b$, wyznaczając sumę $s = a \oplus b$ oraz przeniesienie $c = a \cdot b$; nie posiada wejścia przeniesienia z niższego rzędu.
    \item Sumator pełny (Full Adder) -- dodaje trzy bity $a$, $b$ oraz $c_{in}$, wyznaczając sumę oraz przeniesienie do rzędu wyższego; stanowi podstawowy element układów sumujących wielobitowych.
    \item Sumator kaskadowy (Ripple-Carry Adder) -- powstaje przez połączenie kaskadowe sumatorów pełnych, w którym przeniesienie z każdego rzędu podawane jest na wejście przeniesienia rzędu następnego.
\end{itemize}

W praktycznych układach wyróżnia się również bardziej zaawansowane struktury, takie jak sumatory z przyspieszonym przeniesieniem (Carry-Lookahead Adder), w których przeniesienia wyznaczane są równolegle dla wszystkich pozycji, co pozwala zredukować czas propagacji sygnału kosztem zwiększenia złożoności układu.

Ze względu na liczbę bitów obsługiwanych przez pojedynczy stopień można wyróżnić sumatory jedno- oraz wielobitowe. Zwiększenie liczby bitów powoduje wzrost złożoności układu, dlatego istotnym elementem projektowania układów cyfrowych jest ich odpowiednia minimalizacja.

W niniejszym projekcie zastosowany zostanie pojedynczy sumator pełny. Układ będzie dodawał trzy bity wejściowe $a$, $b$ oraz $c_{in}$, a jego wyjścia $s$ oraz $c_{out}$ będą informowały o wartości sumy oraz przeniesienia. Następnie układ zostanie poddany procesowi minimalizacji w celu ograniczenia liczby wymaganych bramek logicznych.

% ============================================================
% Punkt 2: Opis techniczny projektu
% Projekt układu sumatora pełnego (Full Adder)
% z minimalizacją bramek logicznych
% ------------------------------------------------------------
% Plik przeznaczony do włączenia w dokumencie głównym Overleaf:
%   % ============================================================
% Punkt 2: Opis techniczny projektu
% Projekt układu sumatora pełnego (Full Adder)
% z minimalizacją bramek logicznych
% ------------------------------------------------------------
% Plik przeznaczony do włączenia w dokumencie głównym Overleaf:
%   % ============================================================
% Punkt 2: Opis techniczny projektu
% Projekt układu sumatora pełnego (Full Adder)
% z minimalizacją bramek logicznych
% ------------------------------------------------------------
% Plik przeznaczony do włączenia w dokumencie głównym Overleaf:
%   \input{02_opis_techniczny}
% Wymagane pakiety w dokumencie głównym:
%   \usepackage[utf8]{inputenc}  \usepackage[T1]{polski}
%   \usepackage{amsmath, amssymb, array, booktabs}
% ============================================================

\section{Opis techniczny projektu}

\subsection{Założenia projektowe}

Celem projektu jest zaprojektowanie układu kombinacyjnego realizującego funkcję sumatora pełnego, a następnie zminimalizowanie liczby wykorzystanych bramek logicznych.

Projektowany układ będzie dodawał trzy bity wejściowe: dwa bity argumentów oraz bit przeniesienia z niższego rzędu:
\begin{equation*}
a, \qquad b, \qquad c_{in},
\end{equation*}
gdzie $c_{in}$ oznacza przeniesienie generowane przez poprzedni (młodszy) stopień kaskady.

Układ posiada trzy wejścia binarne oraz dwa wyjścia. Wyjścia będą informowały o wartości bitu sumy oraz o przeniesieniu do rzędu wyższego.

Przyjęto następujące założenia:
\begin{itemize}
    \item układ dodaje trzy bity wejściowe: $a$, $b$ oraz $c_{in}$,
    \item wszystkie wejścia przyjmują wartości logiczne 0 lub 1,
    \item układ jest układem kombinacyjnym,
    \item układ posiada dwa wyjścia: bit sumy $s$ oraz przeniesienie $c_{out}$,
    \item stan wyjścia $s = 1$ oznacza, że suma wejść jest nieparzysta,
    \item stan wyjścia $c_{out} = 1$ oznacza, że $a + b + c_{in} \geq 2$,
    \item podstawowym celem projektu jest minimalizacja liczby bramek logicznych,
    \item poprawność działania zostanie sprawdzona za pomocą symulacji.
\end{itemize}

Minimalizacja układu ma na celu zmniejszenie jego złożoności. Mniejsza liczba zastosowanych bramek oznacza prostszy schemat logiczny, mniejsze zapotrzebowanie na elementy oraz potencjalnie mniejsze opóźnienia propagacji sygnału.

\subsection{Specyfikacja wejść i wyjść}

Projektowany sumator pełny posiada trzy wejścia:
\begin{itemize}
    \item $a$ -- pierwszy bit argumentu,
    \item $b$ -- drugi bit argumentu,
    \item $c_{in}$ -- przeniesienie z niższego rzędu.
\end{itemize}

Wyjściami układu są:
\begin{itemize}
    \item $s$ -- bit sumy,
    \item $c_{out}$ -- przeniesienie do rzędu wyższego.
\end{itemize}

Zależność pomiędzy wejściami i wyjściami można zapisać jako:
\begin{equation*}
s, \, c_{out} = f(a, b, c_{in}).
\end{equation*}

Przykładowo, dla:
\begin{equation*}
a = 1, \qquad b = 1, \qquad c_{in} = 0,
\end{equation*}
otrzymujemy:
\begin{equation*}
s = 0, \qquad c_{out} = 1.
\end{equation*}

Natomiast dla:
\begin{equation*}
a = 1, \qquad b = 1, \qquad c_{in} = 1,
\end{equation*}
otrzymujemy:
\begin{equation*}
s = 1, \qquad c_{out} = 1.
\end{equation*}

\vspace{0.5em}
\noindent Tabela sygnałów:

\begin{center}
\begin{tabular}{|c|c|l|}
\hline
\textbf{Sygnał} & \textbf{Rodzaj} & \textbf{Opis} \\
\hline
$a$       & wejście & pierwszy bit argumentu \\
$b$       & wejście & drugi bit argumentu \\
$c_{in}$  & wejście & przeniesienie z niższego rzędu \\
$s$       & wyjście & bit sumy \\
$c_{out}$ & wyjście & przeniesienie do rzędu wyższego \\
\hline
\end{tabular}
\end{center}

\subsection{Analiza logiczna układu}

Sumator pełny wykonuje arytmetyczne dodawanie trzech bitów, w wyniku czego może powstać wartość z zakresu od 0 do 3, wymagająca dwóch bitów zapisu: $c_{out}s$.

Aby suma była równa 1, nieparzysta liczba sygnałów wejściowych musi przyjmować wartość wysoką. Warunek ten realizuje dwuargumentowa operacja alternatywy rozłącznej, zastosowana dwukrotnie:
\begin{equation*}
s = a \oplus b \oplus c_{in}.
\end{equation*}

Przeniesienie pojawia się wtedy, gdy co najmniej dwa z trzech wejść znajdują się w stanie wysokim. Funkcję wyjściową przeniesienia można zapisać jako:
\begin{equation*}
c_{out} = a \cdot b \vee a \cdot c_{in} \vee b \cdot c_{in}.
\end{equation*}

Jest to tzw. funkcja większości, która przyjmuje wartość 1, gdy większość (co najmniej dwa) sygnałów wejściowych ma wartość 1.

Wprowadzając oznaczenie pomocnicze:
\begin{equation*}
c_1 = a \oplus b,
\end{equation*}
można obydwie funkcje zapisać w postaci:
\begin{equation*}
s = c_1 \oplus c_{in}, \qquad c_{out} = a \cdot b \vee c_1 \cdot c_{in}.
\end{equation*}

Jeżeli nie korzystamy z bramek XOR, alternatywę rozłączną można zapisać przy pomocy podstawowych operacji AND, OR i NOT:
\begin{equation*}
a \oplus b = a \cdot b' \vee a' \cdot b.
\end{equation*}

Stąd:
\begin{equation*}
s = a \cdot b' \cdot c_{in}' \vee a \cdot b \cdot c_{in} \vee a' \cdot b \cdot c_{in}' \vee a' \cdot b' \cdot c_{in}.
\end{equation*}

Ta postać będzie szczególnie istotna podczas dalszej analizy i minimalizacji układu.

\subsection{Tablica prawdy}

Ponieważ układ posiada trzy wejścia binarne, liczba możliwych kombinacji wynosi:
\begin{equation*}
2^3 = 8.
\end{equation*}

Pełna tablica prawdy przedstawia się następująco:

\begin{center}
\begin{tabular}{|c|c|c|c|c|c|}
\hline
$a$ & $b$ & $c_{in}$ & $a+b+c_{in}$ & $s$ & $c_{out}$ \\
\hline
0 & 0 & 0 & 0 & 0 & 0 \\
0 & 0 & 1 & 1 & 1 & 0 \\
0 & 1 & 0 & 1 & 1 & 0 \\
0 & 1 & 1 & 2 & 0 & 1 \\
1 & 0 & 0 & 1 & 1 & 0 \\
1 & 0 & 1 & 2 & 0 & 1 \\
1 & 1 & 0 & 2 & 0 & 1 \\
1 & 1 & 1 & 3 & 1 & 1 \\
\hline
\end{tabular}
\end{center}

Wartość wyjścia $s = 1$ występuje wyłącznie dla kombinacji:
\begin{equation*}
(0,0,1), \quad (0,1,0), \quad (1,0,0), \quad (1,1,1).
\end{equation*}

Funkcję w postaci sumy iloczynów mintermów można więc zapisać jako:
\begin{equation*}
s = \Sigma m(1, 2, 4, 7).
\end{equation*}

Po rozpisaniu:
\begin{equation*}
s = a' \cdot b' \cdot c_{in} \vee a' \cdot b \cdot c_{in}' \vee a \cdot b' \cdot c_{in}' \vee a \cdot b \cdot c_{in}.
\end{equation*}

Analogicznie dla wyjścia przeniesienia, gdzie $c_{out} = 1$ dla kombinacji $(0,1,1)$, $(1,0,1)$, $(1,1,0)$ oraz $(1,1,1)$:
\begin{equation*}
c_{out} = \Sigma m(3, 5, 6, 7),
\end{equation*}
\begin{equation*}
c_{out} = a' \cdot b \cdot c_{in} \vee a \cdot b' \cdot c_{in} \vee a \cdot b \cdot c_{in}' \vee a \cdot b \cdot c_{in}.
\end{equation*}

Jest to postać kanoniczna, która będzie punktem wyjścia do minimalizacji.

\subsection{Mapa Karnaugha}

Do uproszczenia funkcji wykorzystujemy mapy Karnaugha dla trzech zmiennych, osobno dla każdego z wyjść.

Przyjmujemy:
\begin{itemize}
    \item wiersze: pary zmiennych $ab$,
    \item kolumny: zmienna $c_{in}$,
    \item wiersze uporządkowane w kolejności kodu Graya: $00, 01, 11, 10$.
\end{itemize}

Mapa Karnaugha dla wyjścia $s$ przyjmuje postać:

\begin{center}
\begin{tabular}{|c|c|c|}
\hline
$ab \backslash c_{in}$ & \textbf{0} & \textbf{1} \\
\hline
\textbf{00} & 0 & 1 \\
\hline
\textbf{01} & 1 & 0 \\
\hline
\textbf{11} & 0 & 1 \\
\hline
\textbf{10} & 1 & 0 \\
\hline
\end{tabular}
\end{center}

Jedynki znajdują się więc na przekątnej mapy. W tym przypadku nie można połączyć sąsiednich jedynek w większe grupy, ponieważ żadna z nich nie posiada sąsiedniej jedynki zgodnie z zasadami map Karnaugha. Oznacza to, że bez wykorzystania dodatkowych zależności funkcja $s$ pozostaje w postaci czterech iloczynów trójzmiennych. Jednocześnie sama struktura problemu pozwala zauważyć, że funkcję można znacznie efektywniej zrealizować, rozpoznając w niej dwukrotne zastosowanie alternatywy rozłącznej (XOR).

Mapa Karnaugha dla wyjścia $c_{out}$ przyjmuje postać:

\begin{center}
\begin{tabular}{|c|c|c|}
\hline
$ab \backslash c_{in}$ & \textbf{0} & \textbf{1} \\
\hline
\textbf{00} & 0 & 0 \\
\hline
\textbf{01} & 0 & 1 \\
\hline
\textbf{11} & 1 & 1 \\
\hline
\textbf{10} & 0 & 1 \\
\hline
\end{tabular}
\end{center}

Dla funkcji $c_{out}$ możliwe jest utworzenie trzech grup po dwie jedynki: pary w wierszu $ab = 11$, pary w kolumnie $c_{in} = 1$ obejmującej wiersze $01$ i $11$ oraz pary w kolumnie $c_{in} = 1$ obejmującej wiersze $11$ i $10$. W wyniku grupowania otrzymujemy zredukowaną postać funkcji większości:
\begin{equation*}
c_{out} = a \cdot b \vee b \cdot c_{in} \vee a \cdot c_{in}.
\end{equation*}

\subsection{Minimalizacja funkcji logicznej}

Dla wyjścia sumy, na podstawie analizy struktury funkcji, otrzymujemy:
\begin{equation*}
s = a \oplus b \oplus c_{in}.
\end{equation*}

Wprowadzając sygnał pomocniczy:
\begin{equation*}
c_1 = a \oplus b,
\end{equation*}
zapis ten upraszcza się do:
\begin{equation*}
s = c_1 \oplus c_{in}.
\end{equation*}

Dla wyjścia przeniesienia, w wyniku grupowania na mapie Karnaugha:
\begin{equation*}
c_{out} = a \cdot b \vee a \cdot c_{in} \vee b \cdot c_{in}.
\end{equation*}

Alternatywnie, wykorzystując wspólny sygnał $c_1$, przeniesienie można zapisać jako:
\begin{equation*}
c_{out} = a \cdot b \vee c_1 \cdot c_{in}.
\end{equation*}

Jest to postać znacznie bardziej przejrzysta od postaci kanonicznej.

Jeżeli w realizacji można zastosować bramki XOR, każdy z członów może zostać zrealizowany pojedynczą bramką:
\begin{equation*}
c_1 = a \oplus b, \qquad s = c_1 \oplus c_{in}.
\end{equation*}

Ostateczna realizacja wymaga więc:
\begin{itemize}
    \item 2 $\times$ XOR ($c_1$ oraz $s$),
    \item 2 $\times$ AND2 ($a \cdot b$ oraz $c_1 \cdot c_{in}$),
    \item 1 $\times$ OR2 (sumowanie członów przeniesienia).
\end{itemize}

Czyli łącznie 5 bramek.

To rozwiązanie będzie naszym wariantem zoptymalizowanym pod względem liczby bramek, o ile w dalszej części przyjmiemy XOR jako dostępny element logiczny.

\subsection{Porównanie liczby bramek przed i po minimalizacji}

Dla realizacji funkcji w sposób bezpośredni, na podstawie postaci kanonicznej, konieczne byłoby zastosowanie wielu bramek AND, NOT oraz OR.

Postać kanoniczna wyjścia sumy:
\begin{equation*}
s = a' \cdot b' \cdot c_{in} \vee a' \cdot b \cdot c_{in}' \vee a \cdot b' \cdot c_{in}' \vee a \cdot b \cdot c_{in}
\end{equation*}
wymaga utworzenia czterech iloczynów trójzmiennych, a następnie ich zsumowania. Realizacja ta wymaga: 3 bramek NOT (negacje $a$, $b$, $c_{in}$), 4 bramek AND3 oraz 1 bramki OR4, czyli łącznie 8 bramek. Analogicznie dla wyjścia $c_{out}$ kolejne 8 bramek, przy czym bramki NOT mogą być wspólne.

Łącznie realizacja bezpośrednia wymaga około 13--16 bramek logicznych o głębokości trzech poziomów.

Po zastosowaniu zależności wynikającej z alternatywy rozłącznej oraz grupowania na mapach Karnaugha otrzymujemy:
\begin{equation*}
s = a \oplus b \oplus c_{in}, \qquad c_{out} = a \cdot b \vee c_1 \cdot c_{in}, \qquad c_1 = a \oplus b.
\end{equation*}

W rezultacie układ można zrealizować przy użyciu dwóch bramek XOR, dwóch bramek AND oraz jednej bramki OR, czyli łącznie 5 bramek.

Porównanie obu wariantów przedstawia tabela:

\begin{center}
\begin{tabular}{|l|c|c|}
\hline
\textbf{Wariant realizacji} & \textbf{Liczba bramek} & \textbf{Głębokość logiczna} \\
\hline
Postać kanoniczna (SOP) & ok. 13--16 & 3 poziomy \\
Wariant zminimalizowany & 5 & 3 poziomy \\
\hline
\end{tabular}
\end{center}

Zastosowanie wspólnego sygnału $c_1 = a \oplus b$ pozwala dodatkowo wyeliminować powtarzanie tej samej operacji dla obu wyjść, co stanowi dodatkową oszczędność w porównaniu z realizacją, w której każda funkcja minimalizowana jest niezależnie.

% Wymagane pakiety w dokumencie głównym:
%   \usepackage[utf8]{inputenc}  \usepackage[T1]{polski}
%   \usepackage{amsmath, amssymb, array, booktabs}
% ============================================================

\section{Opis techniczny projektu}

\subsection{Założenia projektowe}

Celem projektu jest zaprojektowanie układu kombinacyjnego realizującego funkcję sumatora pełnego, a następnie zminimalizowanie liczby wykorzystanych bramek logicznych.

Projektowany układ będzie dodawał trzy bity wejściowe: dwa bity argumentów oraz bit przeniesienia z niższego rzędu:
\begin{equation*}
a, \qquad b, \qquad c_{in},
\end{equation*}
gdzie $c_{in}$ oznacza przeniesienie generowane przez poprzedni (młodszy) stopień kaskady.

Układ posiada trzy wejścia binarne oraz dwa wyjścia. Wyjścia będą informowały o wartości bitu sumy oraz o przeniesieniu do rzędu wyższego.

Przyjęto następujące założenia:
\begin{itemize}
    \item układ dodaje trzy bity wejściowe: $a$, $b$ oraz $c_{in}$,
    \item wszystkie wejścia przyjmują wartości logiczne 0 lub 1,
    \item układ jest układem kombinacyjnym,
    \item układ posiada dwa wyjścia: bit sumy $s$ oraz przeniesienie $c_{out}$,
    \item stan wyjścia $s = 1$ oznacza, że suma wejść jest nieparzysta,
    \item stan wyjścia $c_{out} = 1$ oznacza, że $a + b + c_{in} \geq 2$,
    \item podstawowym celem projektu jest minimalizacja liczby bramek logicznych,
    \item poprawność działania zostanie sprawdzona za pomocą symulacji.
\end{itemize}

Minimalizacja układu ma na celu zmniejszenie jego złożoności. Mniejsza liczba zastosowanych bramek oznacza prostszy schemat logiczny, mniejsze zapotrzebowanie na elementy oraz potencjalnie mniejsze opóźnienia propagacji sygnału.

\subsection{Specyfikacja wejść i wyjść}

Projektowany sumator pełny posiada trzy wejścia:
\begin{itemize}
    \item $a$ -- pierwszy bit argumentu,
    \item $b$ -- drugi bit argumentu,
    \item $c_{in}$ -- przeniesienie z niższego rzędu.
\end{itemize}

Wyjściami układu są:
\begin{itemize}
    \item $s$ -- bit sumy,
    \item $c_{out}$ -- przeniesienie do rzędu wyższego.
\end{itemize}

Zależność pomiędzy wejściami i wyjściami można zapisać jako:
\begin{equation*}
s, \, c_{out} = f(a, b, c_{in}).
\end{equation*}

Przykładowo, dla:
\begin{equation*}
a = 1, \qquad b = 1, \qquad c_{in} = 0,
\end{equation*}
otrzymujemy:
\begin{equation*}
s = 0, \qquad c_{out} = 1.
\end{equation*}

Natomiast dla:
\begin{equation*}
a = 1, \qquad b = 1, \qquad c_{in} = 1,
\end{equation*}
otrzymujemy:
\begin{equation*}
s = 1, \qquad c_{out} = 1.
\end{equation*}

\vspace{0.5em}
\noindent Tabela sygnałów:

\begin{center}
\begin{tabular}{|c|c|l|}
\hline
\textbf{Sygnał} & \textbf{Rodzaj} & \textbf{Opis} \\
\hline
$a$       & wejście & pierwszy bit argumentu \\
$b$       & wejście & drugi bit argumentu \\
$c_{in}$  & wejście & przeniesienie z niższego rzędu \\
$s$       & wyjście & bit sumy \\
$c_{out}$ & wyjście & przeniesienie do rzędu wyższego \\
\hline
\end{tabular}
\end{center}

\subsection{Analiza logiczna układu}

Sumator pełny wykonuje arytmetyczne dodawanie trzech bitów, w wyniku czego może powstać wartość z zakresu od 0 do 3, wymagająca dwóch bitów zapisu: $c_{out}s$.

Aby suma była równa 1, nieparzysta liczba sygnałów wejściowych musi przyjmować wartość wysoką. Warunek ten realizuje dwuargumentowa operacja alternatywy rozłącznej, zastosowana dwukrotnie:
\begin{equation*}
s = a \oplus b \oplus c_{in}.
\end{equation*}

Przeniesienie pojawia się wtedy, gdy co najmniej dwa z trzech wejść znajdują się w stanie wysokim. Funkcję wyjściową przeniesienia można zapisać jako:
\begin{equation*}
c_{out} = a \cdot b \vee a \cdot c_{in} \vee b \cdot c_{in}.
\end{equation*}

Jest to tzw. funkcja większości, która przyjmuje wartość 1, gdy większość (co najmniej dwa) sygnałów wejściowych ma wartość 1.

Wprowadzając oznaczenie pomocnicze:
\begin{equation*}
c_1 = a \oplus b,
\end{equation*}
można obydwie funkcje zapisać w postaci:
\begin{equation*}
s = c_1 \oplus c_{in}, \qquad c_{out} = a \cdot b \vee c_1 \cdot c_{in}.
\end{equation*}

Jeżeli nie korzystamy z bramek XOR, alternatywę rozłączną można zapisać przy pomocy podstawowych operacji AND, OR i NOT:
\begin{equation*}
a \oplus b = a \cdot b' \vee a' \cdot b.
\end{equation*}

Stąd:
\begin{equation*}
s = a \cdot b' \cdot c_{in}' \vee a \cdot b \cdot c_{in} \vee a' \cdot b \cdot c_{in}' \vee a' \cdot b' \cdot c_{in}.
\end{equation*}

Ta postać będzie szczególnie istotna podczas dalszej analizy i minimalizacji układu.

\subsection{Tablica prawdy}

Ponieważ układ posiada trzy wejścia binarne, liczba możliwych kombinacji wynosi:
\begin{equation*}
2^3 = 8.
\end{equation*}

Pełna tablica prawdy przedstawia się następująco:

\begin{center}
\begin{tabular}{|c|c|c|c|c|c|}
\hline
$a$ & $b$ & $c_{in}$ & $a+b+c_{in}$ & $s$ & $c_{out}$ \\
\hline
0 & 0 & 0 & 0 & 0 & 0 \\
0 & 0 & 1 & 1 & 1 & 0 \\
0 & 1 & 0 & 1 & 1 & 0 \\
0 & 1 & 1 & 2 & 0 & 1 \\
1 & 0 & 0 & 1 & 1 & 0 \\
1 & 0 & 1 & 2 & 0 & 1 \\
1 & 1 & 0 & 2 & 0 & 1 \\
1 & 1 & 1 & 3 & 1 & 1 \\
\hline
\end{tabular}
\end{center}

Wartość wyjścia $s = 1$ występuje wyłącznie dla kombinacji:
\begin{equation*}
(0,0,1), \quad (0,1,0), \quad (1,0,0), \quad (1,1,1).
\end{equation*}

Funkcję w postaci sumy iloczynów mintermów można więc zapisać jako:
\begin{equation*}
s = \Sigma m(1, 2, 4, 7).
\end{equation*}

Po rozpisaniu:
\begin{equation*}
s = a' \cdot b' \cdot c_{in} \vee a' \cdot b \cdot c_{in}' \vee a \cdot b' \cdot c_{in}' \vee a \cdot b \cdot c_{in}.
\end{equation*}

Analogicznie dla wyjścia przeniesienia, gdzie $c_{out} = 1$ dla kombinacji $(0,1,1)$, $(1,0,1)$, $(1,1,0)$ oraz $(1,1,1)$:
\begin{equation*}
c_{out} = \Sigma m(3, 5, 6, 7),
\end{equation*}
\begin{equation*}
c_{out} = a' \cdot b \cdot c_{in} \vee a \cdot b' \cdot c_{in} \vee a \cdot b \cdot c_{in}' \vee a \cdot b \cdot c_{in}.
\end{equation*}

Jest to postać kanoniczna, która będzie punktem wyjścia do minimalizacji.

\subsection{Mapa Karnaugha}

Do uproszczenia funkcji wykorzystujemy mapy Karnaugha dla trzech zmiennych, osobno dla każdego z wyjść.

Przyjmujemy:
\begin{itemize}
    \item wiersze: pary zmiennych $ab$,
    \item kolumny: zmienna $c_{in}$,
    \item wiersze uporządkowane w kolejności kodu Graya: $00, 01, 11, 10$.
\end{itemize}

Mapa Karnaugha dla wyjścia $s$ przyjmuje postać:

\begin{center}
\begin{tabular}{|c|c|c|}
\hline
$ab \backslash c_{in}$ & \textbf{0} & \textbf{1} \\
\hline
\textbf{00} & 0 & 1 \\
\hline
\textbf{01} & 1 & 0 \\
\hline
\textbf{11} & 0 & 1 \\
\hline
\textbf{10} & 1 & 0 \\
\hline
\end{tabular}
\end{center}

Jedynki znajdują się więc na przekątnej mapy. W tym przypadku nie można połączyć sąsiednich jedynek w większe grupy, ponieważ żadna z nich nie posiada sąsiedniej jedynki zgodnie z zasadami map Karnaugha. Oznacza to, że bez wykorzystania dodatkowych zależności funkcja $s$ pozostaje w postaci czterech iloczynów trójzmiennych. Jednocześnie sama struktura problemu pozwala zauważyć, że funkcję można znacznie efektywniej zrealizować, rozpoznając w niej dwukrotne zastosowanie alternatywy rozłącznej (XOR).

Mapa Karnaugha dla wyjścia $c_{out}$ przyjmuje postać:

\begin{center}
\begin{tabular}{|c|c|c|}
\hline
$ab \backslash c_{in}$ & \textbf{0} & \textbf{1} \\
\hline
\textbf{00} & 0 & 0 \\
\hline
\textbf{01} & 0 & 1 \\
\hline
\textbf{11} & 1 & 1 \\
\hline
\textbf{10} & 0 & 1 \\
\hline
\end{tabular}
\end{center}

Dla funkcji $c_{out}$ możliwe jest utworzenie trzech grup po dwie jedynki: pary w wierszu $ab = 11$, pary w kolumnie $c_{in} = 1$ obejmującej wiersze $01$ i $11$ oraz pary w kolumnie $c_{in} = 1$ obejmującej wiersze $11$ i $10$. W wyniku grupowania otrzymujemy zredukowaną postać funkcji większości:
\begin{equation*}
c_{out} = a \cdot b \vee b \cdot c_{in} \vee a \cdot c_{in}.
\end{equation*}

\subsection{Minimalizacja funkcji logicznej}

Dla wyjścia sumy, na podstawie analizy struktury funkcji, otrzymujemy:
\begin{equation*}
s = a \oplus b \oplus c_{in}.
\end{equation*}

Wprowadzając sygnał pomocniczy:
\begin{equation*}
c_1 = a \oplus b,
\end{equation*}
zapis ten upraszcza się do:
\begin{equation*}
s = c_1 \oplus c_{in}.
\end{equation*}

Dla wyjścia przeniesienia, w wyniku grupowania na mapie Karnaugha:
\begin{equation*}
c_{out} = a \cdot b \vee a \cdot c_{in} \vee b \cdot c_{in}.
\end{equation*}

Alternatywnie, wykorzystując wspólny sygnał $c_1$, przeniesienie można zapisać jako:
\begin{equation*}
c_{out} = a \cdot b \vee c_1 \cdot c_{in}.
\end{equation*}

Jest to postać znacznie bardziej przejrzysta od postaci kanonicznej.

Jeżeli w realizacji można zastosować bramki XOR, każdy z członów może zostać zrealizowany pojedynczą bramką:
\begin{equation*}
c_1 = a \oplus b, \qquad s = c_1 \oplus c_{in}.
\end{equation*}

Ostateczna realizacja wymaga więc:
\begin{itemize}
    \item 2 $\times$ XOR ($c_1$ oraz $s$),
    \item 2 $\times$ AND2 ($a \cdot b$ oraz $c_1 \cdot c_{in}$),
    \item 1 $\times$ OR2 (sumowanie członów przeniesienia).
\end{itemize}

Czyli łącznie 5 bramek.

To rozwiązanie będzie naszym wariantem zoptymalizowanym pod względem liczby bramek, o ile w dalszej części przyjmiemy XOR jako dostępny element logiczny.

\subsection{Porównanie liczby bramek przed i po minimalizacji}

Dla realizacji funkcji w sposób bezpośredni, na podstawie postaci kanonicznej, konieczne byłoby zastosowanie wielu bramek AND, NOT oraz OR.

Postać kanoniczna wyjścia sumy:
\begin{equation*}
s = a' \cdot b' \cdot c_{in} \vee a' \cdot b \cdot c_{in}' \vee a \cdot b' \cdot c_{in}' \vee a \cdot b \cdot c_{in}
\end{equation*}
wymaga utworzenia czterech iloczynów trójzmiennych, a następnie ich zsumowania. Realizacja ta wymaga: 3 bramek NOT (negacje $a$, $b$, $c_{in}$), 4 bramek AND3 oraz 1 bramki OR4, czyli łącznie 8 bramek. Analogicznie dla wyjścia $c_{out}$ kolejne 8 bramek, przy czym bramki NOT mogą być wspólne.

Łącznie realizacja bezpośrednia wymaga około 13--16 bramek logicznych o głębokości trzech poziomów.

Po zastosowaniu zależności wynikającej z alternatywy rozłącznej oraz grupowania na mapach Karnaugha otrzymujemy:
\begin{equation*}
s = a \oplus b \oplus c_{in}, \qquad c_{out} = a \cdot b \vee c_1 \cdot c_{in}, \qquad c_1 = a \oplus b.
\end{equation*}

W rezultacie układ można zrealizować przy użyciu dwóch bramek XOR, dwóch bramek AND oraz jednej bramki OR, czyli łącznie 5 bramek.

Porównanie obu wariantów przedstawia tabela:

\begin{center}
\begin{tabular}{|l|c|c|}
\hline
\textbf{Wariant realizacji} & \textbf{Liczba bramek} & \textbf{Głębokość logiczna} \\
\hline
Postać kanoniczna (SOP) & ok. 13--16 & 3 poziomy \\
Wariant zminimalizowany & 5 & 3 poziomy \\
\hline
\end{tabular}
\end{center}

Zastosowanie wspólnego sygnału $c_1 = a \oplus b$ pozwala dodatkowo wyeliminować powtarzanie tej samej operacji dla obu wyjść, co stanowi dodatkową oszczędność w porównaniu z realizacją, w której każda funkcja minimalizowana jest niezależnie.

% Wymagane pakiety w dokumencie głównym:
%   \usepackage[utf8]{inputenc}  \usepackage[T1]{polski}
%   \usepackage{amsmath, amssymb, array, booktabs}
% ============================================================

\section{Opis techniczny projektu}

\subsection{Założenia projektowe}

Celem projektu jest zaprojektowanie układu kombinacyjnego realizującego funkcję sumatora pełnego, a następnie zminimalizowanie liczby wykorzystanych bramek logicznych.

Projektowany układ będzie dodawał trzy bity wejściowe: dwa bity argumentów oraz bit przeniesienia z niższego rzędu:
\begin{equation*}
a, \qquad b, \qquad c_{in},
\end{equation*}
gdzie $c_{in}$ oznacza przeniesienie generowane przez poprzedni (młodszy) stopień kaskady.

Układ posiada trzy wejścia binarne oraz dwa wyjścia. Wyjścia będą informowały o wartości bitu sumy oraz o przeniesieniu do rzędu wyższego.

Przyjęto następujące założenia:
\begin{itemize}
    \item układ dodaje trzy bity wejściowe: $a$, $b$ oraz $c_{in}$,
    \item wszystkie wejścia przyjmują wartości logiczne 0 lub 1,
    \item układ jest układem kombinacyjnym,
    \item układ posiada dwa wyjścia: bit sumy $s$ oraz przeniesienie $c_{out}$,
    \item stan wyjścia $s = 1$ oznacza, że suma wejść jest nieparzysta,
    \item stan wyjścia $c_{out} = 1$ oznacza, że $a + b + c_{in} \geq 2$,
    \item podstawowym celem projektu jest minimalizacja liczby bramek logicznych,
    \item poprawność działania zostanie sprawdzona za pomocą symulacji.
\end{itemize}

Minimalizacja układu ma na celu zmniejszenie jego złożoności. Mniejsza liczba zastosowanych bramek oznacza prostszy schemat logiczny, mniejsze zapotrzebowanie na elementy oraz potencjalnie mniejsze opóźnienia propagacji sygnału.

\subsection{Specyfikacja wejść i wyjść}

Projektowany sumator pełny posiada trzy wejścia:
\begin{itemize}
    \item $a$ -- pierwszy bit argumentu,
    \item $b$ -- drugi bit argumentu,
    \item $c_{in}$ -- przeniesienie z niższego rzędu.
\end{itemize}

Wyjściami układu są:
\begin{itemize}
    \item $s$ -- bit sumy,
    \item $c_{out}$ -- przeniesienie do rzędu wyższego.
\end{itemize}

Zależność pomiędzy wejściami i wyjściami można zapisać jako:
\begin{equation*}
s, \, c_{out} = f(a, b, c_{in}).
\end{equation*}

Przykładowo, dla:
\begin{equation*}
a = 1, \qquad b = 1, \qquad c_{in} = 0,
\end{equation*}
otrzymujemy:
\begin{equation*}
s = 0, \qquad c_{out} = 1.
\end{equation*}

Natomiast dla:
\begin{equation*}
a = 1, \qquad b = 1, \qquad c_{in} = 1,
\end{equation*}
otrzymujemy:
\begin{equation*}
s = 1, \qquad c_{out} = 1.
\end{equation*}

\vspace{0.5em}
\noindent Tabela sygnałów:

\begin{center}
\begin{tabular}{|c|c|l|}
\hline
\textbf{Sygnał} & \textbf{Rodzaj} & \textbf{Opis} \\
\hline
$a$       & wejście & pierwszy bit argumentu \\
$b$       & wejście & drugi bit argumentu \\
$c_{in}$  & wejście & przeniesienie z niższego rzędu \\
$s$       & wyjście & bit sumy \\
$c_{out}$ & wyjście & przeniesienie do rzędu wyższego \\
\hline
\end{tabular}
\end{center}

\subsection{Analiza logiczna układu}

Sumator pełny wykonuje arytmetyczne dodawanie trzech bitów, w wyniku czego może powstać wartość z zakresu od 0 do 3, wymagająca dwóch bitów zapisu: $c_{out}s$.

Aby suma była równa 1, nieparzysta liczba sygnałów wejściowych musi przyjmować wartość wysoką. Warunek ten realizuje dwuargumentowa operacja alternatywy rozłącznej, zastosowana dwukrotnie:
\begin{equation*}
s = a \oplus b \oplus c_{in}.
\end{equation*}

Przeniesienie pojawia się wtedy, gdy co najmniej dwa z trzech wejść znajdują się w stanie wysokim. Funkcję wyjściową przeniesienia można zapisać jako:
\begin{equation*}
c_{out} = a \cdot b \vee a \cdot c_{in} \vee b \cdot c_{in}.
\end{equation*}

Jest to tzw. funkcja większości, która przyjmuje wartość 1, gdy większość (co najmniej dwa) sygnałów wejściowych ma wartość 1.

Wprowadzając oznaczenie pomocnicze:
\begin{equation*}
c_1 = a \oplus b,
\end{equation*}
można obydwie funkcje zapisać w postaci:
\begin{equation*}
s = c_1 \oplus c_{in}, \qquad c_{out} = a \cdot b \vee c_1 \cdot c_{in}.
\end{equation*}

Jeżeli nie korzystamy z bramek XOR, alternatywę rozłączną można zapisać przy pomocy podstawowych operacji AND, OR i NOT:
\begin{equation*}
a \oplus b = a \cdot b' \vee a' \cdot b.
\end{equation*}

Stąd:
\begin{equation*}
s = a \cdot b' \cdot c_{in}' \vee a \cdot b \cdot c_{in} \vee a' \cdot b \cdot c_{in}' \vee a' \cdot b' \cdot c_{in}.
\end{equation*}

Ta postać będzie szczególnie istotna podczas dalszej analizy i minimalizacji układu.

\subsection{Tablica prawdy}

Ponieważ układ posiada trzy wejścia binarne, liczba możliwych kombinacji wynosi:
\begin{equation*}
2^3 = 8.
\end{equation*}

Pełna tablica prawdy przedstawia się następująco:

\begin{center}
\begin{tabular}{|c|c|c|c|c|c|}
\hline
$a$ & $b$ & $c_{in}$ & $a+b+c_{in}$ & $s$ & $c_{out}$ \\
\hline
0 & 0 & 0 & 0 & 0 & 0 \\
0 & 0 & 1 & 1 & 1 & 0 \\
0 & 1 & 0 & 1 & 1 & 0 \\
0 & 1 & 1 & 2 & 0 & 1 \\
1 & 0 & 0 & 1 & 1 & 0 \\
1 & 0 & 1 & 2 & 0 & 1 \\
1 & 1 & 0 & 2 & 0 & 1 \\
1 & 1 & 1 & 3 & 1 & 1 \\
\hline
\end{tabular}
\end{center}

Wartość wyjścia $s = 1$ występuje wyłącznie dla kombinacji:
\begin{equation*}
(0,0,1), \quad (0,1,0), \quad (1,0,0), \quad (1,1,1).
\end{equation*}

Funkcję w postaci sumy iloczynów mintermów można więc zapisać jako:
\begin{equation*}
s = \Sigma m(1, 2, 4, 7).
\end{equation*}

Po rozpisaniu:
\begin{equation*}
s = a' \cdot b' \cdot c_{in} \vee a' \cdot b \cdot c_{in}' \vee a \cdot b' \cdot c_{in}' \vee a \cdot b \cdot c_{in}.
\end{equation*}

Analogicznie dla wyjścia przeniesienia, gdzie $c_{out} = 1$ dla kombinacji $(0,1,1)$, $(1,0,1)$, $(1,1,0)$ oraz $(1,1,1)$:
\begin{equation*}
c_{out} = \Sigma m(3, 5, 6, 7),
\end{equation*}
\begin{equation*}
c_{out} = a' \cdot b \cdot c_{in} \vee a \cdot b' \cdot c_{in} \vee a \cdot b \cdot c_{in}' \vee a \cdot b \cdot c_{in}.
\end{equation*}

Jest to postać kanoniczna, która będzie punktem wyjścia do minimalizacji.

\subsection{Mapa Karnaugha}

Do uproszczenia funkcji wykorzystujemy mapy Karnaugha dla trzech zmiennych, osobno dla każdego z wyjść.

Przyjmujemy:
\begin{itemize}
    \item wiersze: pary zmiennych $ab$,
    \item kolumny: zmienna $c_{in}$,
    \item wiersze uporządkowane w kolejności kodu Graya: $00, 01, 11, 10$.
\end{itemize}

Mapa Karnaugha dla wyjścia $s$ przyjmuje postać:

\begin{center}
\begin{tabular}{|c|c|c|}
\hline
$ab \backslash c_{in}$ & \textbf{0} & \textbf{1} \\
\hline
\textbf{00} & 0 & 1 \\
\hline
\textbf{01} & 1 & 0 \\
\hline
\textbf{11} & 0 & 1 \\
\hline
\textbf{10} & 1 & 0 \\
\hline
\end{tabular}
\end{center}

Jedynki znajdują się więc na przekątnej mapy. W tym przypadku nie można połączyć sąsiednich jedynek w większe grupy, ponieważ żadna z nich nie posiada sąsiedniej jedynki zgodnie z zasadami map Karnaugha. Oznacza to, że bez wykorzystania dodatkowych zależności funkcja $s$ pozostaje w postaci czterech iloczynów trójzmiennych. Jednocześnie sama struktura problemu pozwala zauważyć, że funkcję można znacznie efektywniej zrealizować, rozpoznając w niej dwukrotne zastosowanie alternatywy rozłącznej (XOR).

Mapa Karnaugha dla wyjścia $c_{out}$ przyjmuje postać:

\begin{center}
\begin{tabular}{|c|c|c|}
\hline
$ab \backslash c_{in}$ & \textbf{0} & \textbf{1} \\
\hline
\textbf{00} & 0 & 0 \\
\hline
\textbf{01} & 0 & 1 \\
\hline
\textbf{11} & 1 & 1 \\
\hline
\textbf{10} & 0 & 1 \\
\hline
\end{tabular}
\end{center}

Dla funkcji $c_{out}$ możliwe jest utworzenie trzech grup po dwie jedynki: pary w wierszu $ab = 11$, pary w kolumnie $c_{in} = 1$ obejmującej wiersze $01$ i $11$ oraz pary w kolumnie $c_{in} = 1$ obejmującej wiersze $11$ i $10$. W wyniku grupowania otrzymujemy zredukowaną postać funkcji większości:
\begin{equation*}
c_{out} = a \cdot b \vee b \cdot c_{in} \vee a \cdot c_{in}.
\end{equation*}

\subsection{Minimalizacja funkcji logicznej}

Dla wyjścia sumy, na podstawie analizy struktury funkcji, otrzymujemy:
\begin{equation*}
s = a \oplus b \oplus c_{in}.
\end{equation*}

Wprowadzając sygnał pomocniczy:
\begin{equation*}
c_1 = a \oplus b,
\end{equation*}
zapis ten upraszcza się do:
\begin{equation*}
s = c_1 \oplus c_{in}.
\end{equation*}

Dla wyjścia przeniesienia, w wyniku grupowania na mapie Karnaugha:
\begin{equation*}
c_{out} = a \cdot b \vee a \cdot c_{in} \vee b \cdot c_{in}.
\end{equation*}

Alternatywnie, wykorzystując wspólny sygnał $c_1$, przeniesienie można zapisać jako:
\begin{equation*}
c_{out} = a \cdot b \vee c_1 \cdot c_{in}.
\end{equation*}

Jest to postać znacznie bardziej przejrzysta od postaci kanonicznej.

Jeżeli w realizacji można zastosować bramki XOR, każdy z członów może zostać zrealizowany pojedynczą bramką:
\begin{equation*}
c_1 = a \oplus b, \qquad s = c_1 \oplus c_{in}.
\end{equation*}

Ostateczna realizacja wymaga więc:
\begin{itemize}
    \item 2 $\times$ XOR ($c_1$ oraz $s$),
    \item 2 $\times$ AND2 ($a \cdot b$ oraz $c_1 \cdot c_{in}$),
    \item 1 $\times$ OR2 (sumowanie członów przeniesienia).
\end{itemize}

Czyli łącznie 5 bramek.

To rozwiązanie będzie naszym wariantem zoptymalizowanym pod względem liczby bramek, o ile w dalszej części przyjmiemy XOR jako dostępny element logiczny.

\subsection{Porównanie liczby bramek przed i po minimalizacji}

Dla realizacji funkcji w sposób bezpośredni, na podstawie postaci kanonicznej, konieczne byłoby zastosowanie wielu bramek AND, NOT oraz OR.

Postać kanoniczna wyjścia sumy:
\begin{equation*}
s = a' \cdot b' \cdot c_{in} \vee a' \cdot b \cdot c_{in}' \vee a \cdot b' \cdot c_{in}' \vee a \cdot b \cdot c_{in}
\end{equation*}
wymaga utworzenia czterech iloczynów trójzmiennych, a następnie ich zsumowania. Realizacja ta wymaga: 3 bramek NOT (negacje $a$, $b$, $c_{in}$), 4 bramek AND3 oraz 1 bramki OR4, czyli łącznie 8 bramek. Analogicznie dla wyjścia $c_{out}$ kolejne 8 bramek, przy czym bramki NOT mogą być wspólne.

Łącznie realizacja bezpośrednia wymaga około 13--16 bramek logicznych o głębokości trzech poziomów.

Po zastosowaniu zależności wynikającej z alternatywy rozłącznej oraz grupowania na mapach Karnaugha otrzymujemy:
\begin{equation*}
s = a \oplus b \oplus c_{in}, \qquad c_{out} = a \cdot b \vee c_1 \cdot c_{in}, \qquad c_1 = a \oplus b.
\end{equation*}

W rezultacie układ można zrealizować przy użyciu dwóch bramek XOR, dwóch bramek AND oraz jednej bramki OR, czyli łącznie 5 bramek.

Porównanie obu wariantów przedstawia tabela:

\begin{center}
\begin{tabular}{|l|c|c|}
\hline
\textbf{Wariant realizacji} & \textbf{Liczba bramek} & \textbf{Głębokość logiczna} \\
\hline
Postać kanoniczna (SOP) & ok. 13--16 & 3 poziomy \\
Wariant zminimalizowany & 5 & 3 poziomy \\
\hline
\end{tabular}
\end{center}

Zastosowanie wspólnego sygnału $c_1 = a \oplus b$ pozwala dodatkowo wyeliminować powtarzanie tej samej operacji dla obu wyjść, co stanowi dodatkową oszczędność w porównaniu z realizacją, w której każda funkcja minimalizowana jest niezależnie.

% ============================================================
% Punkt 3: Implementacja i realizacja
% Projekt układu sumatora pełnego (Full Adder)
% z minimalizacją bramek logicznych
% ------------------------------------------------------------
% Plik przeznaczony do włączenia w dokumencie głównym Overleaf:
%   % ============================================================
% Punkt 3: Implementacja i realizacja
% Projekt układu sumatora pełnego (Full Adder)
% z minimalizacją bramek logicznych
% ------------------------------------------------------------
% Plik przeznaczony do włączenia w dokumencie głównym Overleaf:
%   % ============================================================
% Punkt 3: Implementacja i realizacja
% Projekt układu sumatora pełnego (Full Adder)
% z minimalizacją bramek logicznych
% ------------------------------------------------------------
% Plik przeznaczony do włączenia w dokumencie głównym Overleaf:
%   \input{03_implementacja}
% Wymagane pakiety w dokumencie głównym:
%   \usepackage[utf8]{inputenc}  \usepackage[T1]{polski}
%   \usepackage{amsmath, amssymb}
% ============================================================

\section{Implementacja i realizacja}

\subsection{Wybór elementów logicznych}

Do realizacji projektu sumatora pełnego (Full Adder) wybrano podstawowe oraz złożone elementy cyfrowe z rodziny układów logicznych CMOS. Podstawowym kryterium wyboru była możliwość minimalizacji liczby struktur bramkowych oraz optymalizacja czasu propagacji sygnału. W projekcie wykorzystano:

\begin{itemize}
    \item Bramki XOR (alternatywy rozłącznej): stanowiące kluczowy element wyznaczający bit sumy. Bramka XOR generuje stan wysoki (1) na wyjściu tylko wtedy, gdy nieparzysta liczba sygnałów wejściowych jest w stanie wysokim.
    \item Bramki AND (koniunkcji): służące do utworzenia członów iloczynowych funkcji przeniesienia, tj. $a \cdot b$ oraz $c_1 \cdot c_{in}$.
    \item Bramkę OR (alternatywy): służącą do zsumowania logicznego członów iloczynowych, czego wynikiem jest sygnał przeniesienia $c_{out}$.
\end{itemize}

\subsection{Projekt układu sumatora}

Projekt obejmuje sumator pełny dodający trzy bity: $a$, $b$ oraz $c_{in}$. Układ posiada trzy wejścia binarne oraz dwa wyjścia cyfrowe: $s$, które przyjmuje wartość logiczną 1, gdy suma wejść jest nieparzysta, oraz $c_{out}$, które przyjmuje wartość logiczną 1, gdy $a + b + c_{in} \geq 2$.

Równania wyjściowe sumatora przed minimalizacją, zapisane w postaci sumy iloczynów (SOP) dla wszystkich stanów aktywnych wyjść, przybrałyby postać rozbudowanych funkcji logicznych składających się z czterech składników mintermów każda:
\begin{equation*}
s = \Sigma m(1, 2, 4, 7), \qquad c_{out} = \Sigma m(3, 5, 6, 7).
\end{equation*}

Zastosowanie bezpośredniego porównania par bitów za pomocą operacji XOR pozwala na zapisanie funkcji w postaci:
\begin{equation*}
s = a \oplus b \oplus c_{in}, \qquad c_{out} = a \cdot b \vee (a \oplus b) \cdot c_{in}.
\end{equation*}

Rozpisując bramkę XOR na podstawowe operacje logiczne (AND, OR, NOT) funkcja sumy przyjmuje postać:
\begin{equation*}
s = a \cdot b' \cdot c_{in}' \vee a \cdot b \cdot c_{in} \vee a' \cdot b \cdot c_{in}' \vee a' \cdot b' \cdot c_{in}.
\end{equation*}

Dzięki takiemu podejściu, zamiast budować rozległą siatkę bramek AND/OR dla ośmiu mintermów i implementować kilkanaście bramek, układ redukuje się do dwóch struktur sumujących oraz jednego elementu spinającego człony przeniesienia.

\subsection{Schemat logiczny układu}

Schemat logiczny zminimalizowanego sumatora opiera się na strukturze dwupoziomowej:

\begin{itemize}
    \item Poziom pierwszy (Porównanie bitów): bramka XOR pracująca nad parami bitów wejściowych. Bramka ta przyjmuje na wejścia bity $a$ i $b$, generując sygnał pośredni $c_1 = a \oplus b$.
    \item Poziom drugi (Wyznaczenie wyjść): sygnał $c_1$ podawany jest jednocześnie na drugą bramkę XOR (wraz z $c_{in}$), której wyjściem jest bit sumy $s = c_1 \oplus c_{in}$, oraz na bramkę AND (wraz z $c_{in}$), generującą człon $c_1 \cdot c_{in}$.
    \item Poziom trzeci (Agregacja wyniku): wyjście bramki AND oraz sygnał z drugiej bramki AND realizującej iloczyn $a \cdot b$ są doprowadzone do dwuwejściowej bramki OR. Sygnał wyjściowy jest bezpośrednim rezultatem operacji OR na sygnałach cząstkowych:
    \begin{equation*}
    c_{out} = a \cdot b \vee c_1 \cdot c_{in}.
    \end{equation*}
\end{itemize}

Przepływ sygnałów w układzie można przedstawić jako:
\begin{equation*}
a, b \xrightarrow{\text{XOR}} c_1 \xrightarrow{\text{XOR z } c_{in}} s, \qquad a \cdot b \xrightarrow{\text{OR}} c_{out} \xleftarrow{\text{OR}} c_1 \cdot c_{in}.
\end{equation*}

Kluczową zaletą tej struktury jest fakt, że sygnał pośredni $c_1$ wykorzystywany jest dwukrotnie -- zarówno w torze sumy, jak i w torze przeniesienia -- co eliminuje konieczność wykonywania tej samej operacji dwukrotnie.

Schemat logiczny zminimalizowanego układu przedstawia Rysunek~\ref{fig:schemat-fa}.

\begin{figure}[h]
    \centering
    \begin{circuitikz}[american]
        % --- bramka XOR 1 (a, b -> c1) ---
        \draw (0,4.2) node[xor port] (xor1) {};
        % --- bramka XOR 2 (c1, cin -> s) ---
        \draw (3.5,3.2) node[xor port] (xor2) {};
        % --- bramka AND 1 (a, b -> ab) ---
        \draw (0,1.6) node[and port] (and1) {};
        % --- bramka AND 2 (c1, cin -> c1cin) ---
        \draw (3.5,1.2) node[and port] (and2) {};
        % --- bramka OR (ab, c1cin -> cout) ---
        \draw (7,1.4) node[or port] (or1) {};
        % --- wejścia a, b ---
        \draw (xor1.in 1) -- ++(-1.2,0) node[left] {$a$} node[circ] {} coordinate (a);
        \draw (and1.in 1) -- (a) node[circ] {};
        \draw (xor1.in 2) -- ++(-1.2,0) node[left] {$b$} node[circ] {} coordinate (b);
        \draw (and1.in 2) -- (b) node[circ] {};
        % --- wejście cin (rozgałęzienie do xor2 i and2) ---
        \draw (xor2.in 2) -- ++(-1.2,0) coordinate (cin1) node[circ] {};
        \draw (and2.in 2) -- ++(-1.2,0) coordinate (cin2) node[circ] {};
        \draw (cin1) -- ++(-0.4,0) node[left] {$c_{in}$};
        \draw (cin1) -- (cin2);
        % --- sygnał c1: wyjście xor1 -> xor2.in1 i and2.in1 ---
        \draw (xor1.out) -- ++(0.35,0) node[circ] {} coordinate (c1);
        \draw (c1) |- (xor2.in 1);
        \draw (c1) |- (and2.in 1);
        % --- wyjście s ---
        \draw (xor2.out) -- ++(0.8,0) node[right] {$s$};
        % --- wyjście cout ---
        \draw (and1.out) -- ++(0.35,0) coordinate (ab) node[circ] {};
        \draw (ab) |- (or1.in 1);
        \draw (and2.out) -- ++(0.35,0) coordinate (cc) node[circ] {};
        \draw (cc) |- (or1.in 2);
        \draw (or1.out) -- ++(0.8,0) node[right] {$c_{out}$};
        % --- opisy sygnałów pośrednich ---
        \node[above] at (c1) {$c_1$};
        \node[above] at (1.9,2.5) {$a \cdot b$};
        \node[above] at (5.0,1.2) {$c_1 \cdot c_{in}$};
    \end{circuitikz}
    \caption{Schemat logiczny zminimalizowanego sumatora pełnego}
    \label{fig:schemat-fa}
\end{figure}

\subsection{Realizacja zminimalizowanego układu}

W praktycznej realizacji fizycznej lub strukturalnej (np. w strukturach reprogramowalnych FPGA lub na makiecie laboratoryjnej), zminimalizowany układ eliminuje potrzebę stosowania rozbudowanych siatek bramek AND/OR dla pojedynczych mintermów.

Przy użyciu standardowych układów scalonych TTL/CMOS, realizacja wymaga:

\begin{itemize}
    \item Jednego układu scalonego zawierającego bramki XOR (np. 74HC86 -- poczwórna bramka XOR 2-wejściowa).
    \item Jednego układu scalonego z bramkami AND (np. 74HC08 -- poczwórna bramka AND 2-wejściowa).
    \item Jednego układu scalonego z bramkami OR (np. 74HC32 -- poczwórna bramka OR 2-wejściowa).
\end{itemize}

Alternatywnie, jeżeli w dostępnym zestawie elementów nie występują bramki XOR, sumę można zrealizować wyłącznie na bramkach NAND, wykorzystując standardową tożsamość logiczną, co jednak zwiększa liczbę wykorzystanych struktur bramkowych.

Ograniczenie liczby połączeń do minimum znacząco redukuje ryzyko powstawania hazardu statycznego i dynamicznego w układzie, skracając jednocześnie łączny czas propagacji sygnału od wejścia do wyjścia do sumy czasów propagacji zaledwie trzech poziomów bramek.

% Wymagane pakiety w dokumencie głównym:
%   \usepackage[utf8]{inputenc}  \usepackage[T1]{polski}
%   \usepackage{amsmath, amssymb}
% ============================================================

\section{Implementacja i realizacja}

\subsection{Wybór elementów logicznych}

Do realizacji projektu sumatora pełnego (Full Adder) wybrano podstawowe oraz złożone elementy cyfrowe z rodziny układów logicznych CMOS. Podstawowym kryterium wyboru była możliwość minimalizacji liczby struktur bramkowych oraz optymalizacja czasu propagacji sygnału. W projekcie wykorzystano:

\begin{itemize}
    \item Bramki XOR (alternatywy rozłącznej): stanowiące kluczowy element wyznaczający bit sumy. Bramka XOR generuje stan wysoki (1) na wyjściu tylko wtedy, gdy nieparzysta liczba sygnałów wejściowych jest w stanie wysokim.
    \item Bramki AND (koniunkcji): służące do utworzenia członów iloczynowych funkcji przeniesienia, tj. $a \cdot b$ oraz $c_1 \cdot c_{in}$.
    \item Bramkę OR (alternatywy): służącą do zsumowania logicznego członów iloczynowych, czego wynikiem jest sygnał przeniesienia $c_{out}$.
\end{itemize}

\subsection{Projekt układu sumatora}

Projekt obejmuje sumator pełny dodający trzy bity: $a$, $b$ oraz $c_{in}$. Układ posiada trzy wejścia binarne oraz dwa wyjścia cyfrowe: $s$, które przyjmuje wartość logiczną 1, gdy suma wejść jest nieparzysta, oraz $c_{out}$, które przyjmuje wartość logiczną 1, gdy $a + b + c_{in} \geq 2$.

Równania wyjściowe sumatora przed minimalizacją, zapisane w postaci sumy iloczynów (SOP) dla wszystkich stanów aktywnych wyjść, przybrałyby postać rozbudowanych funkcji logicznych składających się z czterech składników mintermów każda:
\begin{equation*}
s = \Sigma m(1, 2, 4, 7), \qquad c_{out} = \Sigma m(3, 5, 6, 7).
\end{equation*}

Zastosowanie bezpośredniego porównania par bitów za pomocą operacji XOR pozwala na zapisanie funkcji w postaci:
\begin{equation*}
s = a \oplus b \oplus c_{in}, \qquad c_{out} = a \cdot b \vee (a \oplus b) \cdot c_{in}.
\end{equation*}

Rozpisując bramkę XOR na podstawowe operacje logiczne (AND, OR, NOT) funkcja sumy przyjmuje postać:
\begin{equation*}
s = a \cdot b' \cdot c_{in}' \vee a \cdot b \cdot c_{in} \vee a' \cdot b \cdot c_{in}' \vee a' \cdot b' \cdot c_{in}.
\end{equation*}

Dzięki takiemu podejściu, zamiast budować rozległą siatkę bramek AND/OR dla ośmiu mintermów i implementować kilkanaście bramek, układ redukuje się do dwóch struktur sumujących oraz jednego elementu spinającego człony przeniesienia.

\subsection{Schemat logiczny układu}

Schemat logiczny zminimalizowanego sumatora opiera się na strukturze dwupoziomowej:

\begin{itemize}
    \item Poziom pierwszy (Porównanie bitów): bramka XOR pracująca nad parami bitów wejściowych. Bramka ta przyjmuje na wejścia bity $a$ i $b$, generując sygnał pośredni $c_1 = a \oplus b$.
    \item Poziom drugi (Wyznaczenie wyjść): sygnał $c_1$ podawany jest jednocześnie na drugą bramkę XOR (wraz z $c_{in}$), której wyjściem jest bit sumy $s = c_1 \oplus c_{in}$, oraz na bramkę AND (wraz z $c_{in}$), generującą człon $c_1 \cdot c_{in}$.
    \item Poziom trzeci (Agregacja wyniku): wyjście bramki AND oraz sygnał z drugiej bramki AND realizującej iloczyn $a \cdot b$ są doprowadzone do dwuwejściowej bramki OR. Sygnał wyjściowy jest bezpośrednim rezultatem operacji OR na sygnałach cząstkowych:
    \begin{equation*}
    c_{out} = a \cdot b \vee c_1 \cdot c_{in}.
    \end{equation*}
\end{itemize}

Przepływ sygnałów w układzie można przedstawić jako:
\begin{equation*}
a, b \xrightarrow{\text{XOR}} c_1 \xrightarrow{\text{XOR z } c_{in}} s, \qquad a \cdot b \xrightarrow{\text{OR}} c_{out} \xleftarrow{\text{OR}} c_1 \cdot c_{in}.
\end{equation*}

Kluczową zaletą tej struktury jest fakt, że sygnał pośredni $c_1$ wykorzystywany jest dwukrotnie -- zarówno w torze sumy, jak i w torze przeniesienia -- co eliminuje konieczność wykonywania tej samej operacji dwukrotnie.

Schemat logiczny zminimalizowanego układu przedstawia Rysunek~\ref{fig:schemat-fa}.

\begin{figure}[h]
    \centering
    \begin{circuitikz}[american]
        % --- bramka XOR 1 (a, b -> c1) ---
        \draw (0,4.2) node[xor port] (xor1) {};
        % --- bramka XOR 2 (c1, cin -> s) ---
        \draw (3.5,3.2) node[xor port] (xor2) {};
        % --- bramka AND 1 (a, b -> ab) ---
        \draw (0,1.6) node[and port] (and1) {};
        % --- bramka AND 2 (c1, cin -> c1cin) ---
        \draw (3.5,1.2) node[and port] (and2) {};
        % --- bramka OR (ab, c1cin -> cout) ---
        \draw (7,1.4) node[or port] (or1) {};
        % --- wejścia a, b ---
        \draw (xor1.in 1) -- ++(-1.2,0) node[left] {$a$} node[circ] {} coordinate (a);
        \draw (and1.in 1) -- (a) node[circ] {};
        \draw (xor1.in 2) -- ++(-1.2,0) node[left] {$b$} node[circ] {} coordinate (b);
        \draw (and1.in 2) -- (b) node[circ] {};
        % --- wejście cin (rozgałęzienie do xor2 i and2) ---
        \draw (xor2.in 2) -- ++(-1.2,0) coordinate (cin1) node[circ] {};
        \draw (and2.in 2) -- ++(-1.2,0) coordinate (cin2) node[circ] {};
        \draw (cin1) -- ++(-0.4,0) node[left] {$c_{in}$};
        \draw (cin1) -- (cin2);
        % --- sygnał c1: wyjście xor1 -> xor2.in1 i and2.in1 ---
        \draw (xor1.out) -- ++(0.35,0) node[circ] {} coordinate (c1);
        \draw (c1) |- (xor2.in 1);
        \draw (c1) |- (and2.in 1);
        % --- wyjście s ---
        \draw (xor2.out) -- ++(0.8,0) node[right] {$s$};
        % --- wyjście cout ---
        \draw (and1.out) -- ++(0.35,0) coordinate (ab) node[circ] {};
        \draw (ab) |- (or1.in 1);
        \draw (and2.out) -- ++(0.35,0) coordinate (cc) node[circ] {};
        \draw (cc) |- (or1.in 2);
        \draw (or1.out) -- ++(0.8,0) node[right] {$c_{out}$};
        % --- opisy sygnałów pośrednich ---
        \node[above] at (c1) {$c_1$};
        \node[above] at (1.9,2.5) {$a \cdot b$};
        \node[above] at (5.0,1.2) {$c_1 \cdot c_{in}$};
    \end{circuitikz}
    \caption{Schemat logiczny zminimalizowanego sumatora pełnego}
    \label{fig:schemat-fa}
\end{figure}

\subsection{Realizacja zminimalizowanego układu}

W praktycznej realizacji fizycznej lub strukturalnej (np. w strukturach reprogramowalnych FPGA lub na makiecie laboratoryjnej), zminimalizowany układ eliminuje potrzebę stosowania rozbudowanych siatek bramek AND/OR dla pojedynczych mintermów.

Przy użyciu standardowych układów scalonych TTL/CMOS, realizacja wymaga:

\begin{itemize}
    \item Jednego układu scalonego zawierającego bramki XOR (np. 74HC86 -- poczwórna bramka XOR 2-wejściowa).
    \item Jednego układu scalonego z bramkami AND (np. 74HC08 -- poczwórna bramka AND 2-wejściowa).
    \item Jednego układu scalonego z bramkami OR (np. 74HC32 -- poczwórna bramka OR 2-wejściowa).
\end{itemize}

Alternatywnie, jeżeli w dostępnym zestawie elementów nie występują bramki XOR, sumę można zrealizować wyłącznie na bramkach NAND, wykorzystując standardową tożsamość logiczną, co jednak zwiększa liczbę wykorzystanych struktur bramkowych.

Ograniczenie liczby połączeń do minimum znacząco redukuje ryzyko powstawania hazardu statycznego i dynamicznego w układzie, skracając jednocześnie łączny czas propagacji sygnału od wejścia do wyjścia do sumy czasów propagacji zaledwie trzech poziomów bramek.

% Wymagane pakiety w dokumencie głównym:
%   \usepackage[utf8]{inputenc}  \usepackage[T1]{polski}
%   \usepackage{amsmath, amssymb}
% ============================================================

\section{Implementacja i realizacja}

\subsection{Wybór elementów logicznych}

Do realizacji projektu sumatora pełnego (Full Adder) wybrano podstawowe oraz złożone elementy cyfrowe z rodziny układów logicznych CMOS. Podstawowym kryterium wyboru była możliwość minimalizacji liczby struktur bramkowych oraz optymalizacja czasu propagacji sygnału. W projekcie wykorzystano:

\begin{itemize}
    \item Bramki XOR (alternatywy rozłącznej): stanowiące kluczowy element wyznaczający bit sumy. Bramka XOR generuje stan wysoki (1) na wyjściu tylko wtedy, gdy nieparzysta liczba sygnałów wejściowych jest w stanie wysokim.
    \item Bramki AND (koniunkcji): służące do utworzenia członów iloczynowych funkcji przeniesienia, tj. $a \cdot b$ oraz $c_1 \cdot c_{in}$.
    \item Bramkę OR (alternatywy): służącą do zsumowania logicznego członów iloczynowych, czego wynikiem jest sygnał przeniesienia $c_{out}$.
\end{itemize}

\subsection{Projekt układu sumatora}

Projekt obejmuje sumator pełny dodający trzy bity: $a$, $b$ oraz $c_{in}$. Układ posiada trzy wejścia binarne oraz dwa wyjścia cyfrowe: $s$, które przyjmuje wartość logiczną 1, gdy suma wejść jest nieparzysta, oraz $c_{out}$, które przyjmuje wartość logiczną 1, gdy $a + b + c_{in} \geq 2$.

Równania wyjściowe sumatora przed minimalizacją, zapisane w postaci sumy iloczynów (SOP) dla wszystkich stanów aktywnych wyjść, przybrałyby postać rozbudowanych funkcji logicznych składających się z czterech składników mintermów każda:
\begin{equation*}
s = \Sigma m(1, 2, 4, 7), \qquad c_{out} = \Sigma m(3, 5, 6, 7).
\end{equation*}

Zastosowanie bezpośredniego porównania par bitów za pomocą operacji XOR pozwala na zapisanie funkcji w postaci:
\begin{equation*}
s = a \oplus b \oplus c_{in}, \qquad c_{out} = a \cdot b \vee (a \oplus b) \cdot c_{in}.
\end{equation*}

Rozpisując bramkę XOR na podstawowe operacje logiczne (AND, OR, NOT) funkcja sumy przyjmuje postać:
\begin{equation*}
s = a \cdot b' \cdot c_{in}' \vee a \cdot b \cdot c_{in} \vee a' \cdot b \cdot c_{in}' \vee a' \cdot b' \cdot c_{in}.
\end{equation*}

Dzięki takiemu podejściu, zamiast budować rozległą siatkę bramek AND/OR dla ośmiu mintermów i implementować kilkanaście bramek, układ redukuje się do dwóch struktur sumujących oraz jednego elementu spinającego człony przeniesienia.

\subsection{Schemat logiczny układu}

Schemat logiczny zminimalizowanego sumatora opiera się na strukturze dwupoziomowej:

\begin{itemize}
    \item Poziom pierwszy (Porównanie bitów): bramka XOR pracująca nad parami bitów wejściowych. Bramka ta przyjmuje na wejścia bity $a$ i $b$, generując sygnał pośredni $c_1 = a \oplus b$.
    \item Poziom drugi (Wyznaczenie wyjść): sygnał $c_1$ podawany jest jednocześnie na drugą bramkę XOR (wraz z $c_{in}$), której wyjściem jest bit sumy $s = c_1 \oplus c_{in}$, oraz na bramkę AND (wraz z $c_{in}$), generującą człon $c_1 \cdot c_{in}$.
    \item Poziom trzeci (Agregacja wyniku): wyjście bramki AND oraz sygnał z drugiej bramki AND realizującej iloczyn $a \cdot b$ są doprowadzone do dwuwejściowej bramki OR. Sygnał wyjściowy jest bezpośrednim rezultatem operacji OR na sygnałach cząstkowych:
    \begin{equation*}
    c_{out} = a \cdot b \vee c_1 \cdot c_{in}.
    \end{equation*}
\end{itemize}

Przepływ sygnałów w układzie można przedstawić jako:
\begin{equation*}
a, b \xrightarrow{\text{XOR}} c_1 \xrightarrow{\text{XOR z } c_{in}} s, \qquad a \cdot b \xrightarrow{\text{OR}} c_{out} \xleftarrow{\text{OR}} c_1 \cdot c_{in}.
\end{equation*}

Kluczową zaletą tej struktury jest fakt, że sygnał pośredni $c_1$ wykorzystywany jest dwukrotnie -- zarówno w torze sumy, jak i w torze przeniesienia -- co eliminuje konieczność wykonywania tej samej operacji dwukrotnie.

Schemat logiczny zminimalizowanego układu przedstawia Rysunek~\ref{fig:schemat-fa}.

\begin{figure}[h]
    \centering
    \begin{circuitikz}[american]
        % --- bramka XOR 1 (a, b -> c1) ---
        \draw (0,4.2) node[xor port] (xor1) {};
        % --- bramka XOR 2 (c1, cin -> s) ---
        \draw (3.5,3.2) node[xor port] (xor2) {};
        % --- bramka AND 1 (a, b -> ab) ---
        \draw (0,1.6) node[and port] (and1) {};
        % --- bramka AND 2 (c1, cin -> c1cin) ---
        \draw (3.5,1.2) node[and port] (and2) {};
        % --- bramka OR (ab, c1cin -> cout) ---
        \draw (7,1.4) node[or port] (or1) {};
        % --- wejścia a, b ---
        \draw (xor1.in 1) -- ++(-1.2,0) node[left] {$a$} node[circ] {} coordinate (a);
        \draw (and1.in 1) -- (a) node[circ] {};
        \draw (xor1.in 2) -- ++(-1.2,0) node[left] {$b$} node[circ] {} coordinate (b);
        \draw (and1.in 2) -- (b) node[circ] {};
        % --- wejście cin (rozgałęzienie do xor2 i and2) ---
        \draw (xor2.in 2) -- ++(-1.2,0) coordinate (cin1) node[circ] {};
        \draw (and2.in 2) -- ++(-1.2,0) coordinate (cin2) node[circ] {};
        \draw (cin1) -- ++(-0.4,0) node[left] {$c_{in}$};
        \draw (cin1) -- (cin2);
        % --- sygnał c1: wyjście xor1 -> xor2.in1 i and2.in1 ---
        \draw (xor1.out) -- ++(0.35,0) node[circ] {} coordinate (c1);
        \draw (c1) |- (xor2.in 1);
        \draw (c1) |- (and2.in 1);
        % --- wyjście s ---
        \draw (xor2.out) -- ++(0.8,0) node[right] {$s$};
        % --- wyjście cout ---
        \draw (and1.out) -- ++(0.35,0) coordinate (ab) node[circ] {};
        \draw (ab) |- (or1.in 1);
        \draw (and2.out) -- ++(0.35,0) coordinate (cc) node[circ] {};
        \draw (cc) |- (or1.in 2);
        \draw (or1.out) -- ++(0.8,0) node[right] {$c_{out}$};
        % --- opisy sygnałów pośrednich ---
        \node[above] at (c1) {$c_1$};
        \node[above] at (1.9,2.5) {$a \cdot b$};
        \node[above] at (5.0,1.2) {$c_1 \cdot c_{in}$};
    \end{circuitikz}
    \caption{Schemat logiczny zminimalizowanego sumatora pełnego}
    \label{fig:schemat-fa}
\end{figure}

\subsection{Realizacja zminimalizowanego układu}

W praktycznej realizacji fizycznej lub strukturalnej (np. w strukturach reprogramowalnych FPGA lub na makiecie laboratoryjnej), zminimalizowany układ eliminuje potrzebę stosowania rozbudowanych siatek bramek AND/OR dla pojedynczych mintermów.

Przy użyciu standardowych układów scalonych TTL/CMOS, realizacja wymaga:

\begin{itemize}
    \item Jednego układu scalonego zawierającego bramki XOR (np. 74HC86 -- poczwórna bramka XOR 2-wejściowa).
    \item Jednego układu scalonego z bramkami AND (np. 74HC08 -- poczwórna bramka AND 2-wejściowa).
    \item Jednego układu scalonego z bramkami OR (np. 74HC32 -- poczwórna bramka OR 2-wejściowa).
\end{itemize}

Alternatywnie, jeżeli w dostępnym zestawie elementów nie występują bramki XOR, sumę można zrealizować wyłącznie na bramkach NAND, wykorzystując standardową tożsamość logiczną, co jednak zwiększa liczbę wykorzystanych struktur bramkowych.

Ograniczenie liczby połączeń do minimum znacząco redukuje ryzyko powstawania hazardu statycznego i dynamicznego w układzie, skracając jednocześnie łączny czas propagacji sygnału od wejścia do wyjścia do sumy czasów propagacji zaledwie trzech poziomów bramek.

% ============================================================
% Punkt 4: Symulacja i testowanie
% Projekt układu sumatora pełnego (Full Adder)
% z minimalizacją bramek logicznych
% ------------------------------------------------------------
% Plik przeznaczony do włączenia w dokumencie głównym Overleaf:
%   % ============================================================
% Punkt 4: Symulacja i testowanie
% Projekt układu sumatora pełnego (Full Adder)
% z minimalizacją bramek logicznych
% ------------------------------------------------------------
% Plik przeznaczony do włączenia w dokumencie głównym Overleaf:
%   % ============================================================
% Punkt 4: Symulacja i testowanie
% Projekt układu sumatora pełnego (Full Adder)
% z minimalizacją bramek logicznych
% ------------------------------------------------------------
% Plik przeznaczony do włączenia w dokumencie głównym Overleaf:
%   \input{04_symulacja}
% Wymagane pakiety w dokumencie głównym:
%   \usepackage[utf8]{inputenc}  \usepackage[T1]{polski}
%   \usepackage{amsmath, amssymb, array, booktabs, tikz}
% ============================================================

\section{Symulacja i testowanie}

\subsection{Środowisko symulacyjne}

Weryfikacja poprawności działania zaprojektowanego układu sumatora pełnego została przeprowadzona w środowisku Multisim (lub alternatywnie w programie Logisim-evolution). Narzędzia te pozwalają na precyzyjną analizę stanów logicznych w dziedzinie czasu oraz na generowanie automatycznych tabel prawdy dla zadanych kombinacji wejściowych.

\subsection{Przygotowanie układu do symulacji}

W programie symulacyjnym odwzorowano strukturę opisaną w punkcie 3.3. Do wejść $a$, $b$ oraz $c_{in}$ podłączono cyfrowe generatory stanów (Word Generator) emitujące sekwencje bitowe odpowiadające kolejnym kombinacjom binarnym od 000 do 111. Do wyjść $s$ oraz $c_{out}$ podłączono sondy logiczne oraz analizator stanów logicznych (Logic Analyzer) w celu rejestracji przebiegów czasowych.

\subsection{Scenariusze testowe}

W celu pełnego przetestowania układu zaimplementowano scenariusz obejmujący wszystkie 8 możliwych kombinacji stanów wejściowych $(a, b, c_{in})$. Test podzielono na cztery główne fazy:

\begin{itemize}
    \item Faza 1: Wymuszenie stanu $a = 0$, $b = 0$ i sekwencyjna zmiana sygnału $c_{in}$ od 0 do 1.
    \item Faza 2: Wymuszenie stanu $a = 0$, $b = 1$ i sekwencyjna zmiana sygnału $c_{in}$ od 0 do 1.
    \item Faza 3: Wymuszenie stanu $a = 1$, $b = 0$ i sekwencyjna zmiana sygnału $c_{in}$ od 0 do 1.
    \item Faza 4: Wymuszenie stanu $a = 1$, $b = 1$ i sekwencyjna zmiana sygnału $c_{in}$ od 0 do 1.
\end{itemize}

Szczególna uwaga została zwrócona na punkty, w których następuje wygenerowanie przeniesienia: $a + b + c_{in} \geq 2$.

\subsection{Oczekiwane wyniki symulacji}

Oczekuje się, że sonda logiczna zamontowana na wyjściu $s$ wskaże stan wysoki (1) dokładnie w 4 z 8 przypadków testowych, natomiast sonda na wyjściu $c_{out}$ -- dokładnie w 4 z 8 przypadków. Wykres czasowy z analizatora powinien zobrazować impulsy stanu wysokiego wyłącznie w momentach, gdy suma logiczna wejściowych przebiegów odpowiada wartościom określonym w tablicy prawdy. Dla pozostałych kombinacji wyjścia muszą stabilnie utrzymywać stan niski (0).

Oczekiwane stany wyjść dla wszystkich kombinacji wejściowych przedstawia tabela:

\begin{center}
\begin{tabular}{|c|c|c|c|c|}
\hline
$a$ & $b$ & $c_{in}$ & $s$ & $c_{out}$ \\
\hline
0 & 0 & 0 & 0 & 0 \\
0 & 0 & 1 & 1 & 0 \\
0 & 1 & 0 & 1 & 0 \\
0 & 1 & 1 & 0 & 1 \\
1 & 0 & 0 & 1 & 0 \\
1 & 0 & 1 & 0 & 1 \\
1 & 1 & 0 & 0 & 1 \\
1 & 1 & 1 & 1 & 1 \\
\hline
\end{tabular}
\end{center}

\subsection{Analiza wyników}

Wygenerowany wykres czasowy potwierdził założenia teoretyczne. Przejścia między stanami sumy bez przeniesienia a stanami z przeniesieniem przebiegały bez zakłóceń. Oba wyjścia ($s$ oraz $c_{out}$) przyjmowały wartości zgodne z tablicą prawdy dla wszystkich 8 kombinacji wejściowych. W momentach jednoczesnej zmiany kilku bitów wejściowych (np. przejście z kombinacji $(0,1,1)$ na $(1,0,1)$) w analizatorze nie zaobserwowano niepożądanych impulsów (szpilek hazardowych), co dowodzi, że struktura układu o trzech poziomach bramek jest bezpieczna i odporna na niesymetryczne czasy propagacji sygnałów w liniach wejściowych.

Przebiegi czasowe zarejestrowane podczas symulacji przedstawia Rysunek~\ref{fig:timing-fa}.

\begin{figure}[h]
    \centering
    \begin{tikzpicture}[scale=0.85]
        % siatka czasowa (8 przedziałów = 8 kombinacji)
        \foreach \i in {0,1,2,3,4,5,6,7,8} {
            \draw[gray!40, dashed] (\i*1.6,1.0) -- (\i*1.6,-5.4);
        }
        % etykiety kombinacji (a b cin)
        \node at (0.8,-6.0) {000};
        \node at (2.4,-6.0) {001};
        \node at (4.0,-6.0) {010};
        \node at (5.6,-6.0) {011};
        \node at (7.2,-6.0) {100};
        \node at (8.8,-6.0) {101};
        \node at (10.4,-6.0) {110};
        \node at (12.0,-6.0) {111};
        % a = 00001111
        \node[left] at (-0.2,0.4) {$a$};
        \draw[thick] (0,0) -- (6.4,0) -- (6.4,0.8) -- (12.8,0.8);
        % b = 00110011
        \node[left] at (-0.2,-0.8) {$b$};
        \draw[thick] (0,-1.2) -- (3.2,-1.2) -- (3.2,-0.4) -- (6.4,-0.4) -- (6.4,-1.2) -- (9.6,-1.2) -- (9.6,-0.4) -- (12.8,-0.4);
        % cin = 01010101
        \node[left] at (-0.2,-2.0) {$c_{in}$};
        \draw[thick] (0,-2.4) -- (1.6,-2.4) -- (1.6,-1.6) -- (3.2,-1.6) -- (3.2,-2.4) -- (4.8,-2.4) -- (4.8,-1.6) -- (6.4,-1.6) -- (6.4,-2.4) -- (8.0,-2.4) -- (8.0,-1.6) -- (9.6,-1.6) -- (9.6,-2.4) -- (11.2,-2.4) -- (11.2,-1.6) -- (12.8,-1.6);
        % s = 01101001
        \node[left] at (-0.2,-3.2) {$s$};
        \draw[thick] (0,-3.6) -- (1.6,-3.6) -- (1.6,-2.8) -- (4.8,-2.8) -- (4.8,-3.6) -- (6.4,-3.6) -- (6.4,-2.8) -- (8.0,-2.8) -- (8.0,-3.6) -- (11.2,-3.6) -- (11.2,-2.8) -- (12.8,-2.8);
        % cout = 00010111
        \node[left] at (-0.2,-4.4) {$c_{out}$};
        \draw[thick] (0,-4.8) -- (4.8,-4.8) -- (4.8,-4.0) -- (6.4,-4.0) -- (6.4,-4.8) -- (8.0,-4.8) -- (8.0,-4.0) -- (12.8,-4.0);
        % strzałka czasu
        \draw[->] (0,-5.4) -- (13.2,-5.4) node[right] {$t$};
    \end{tikzpicture}
    \caption{Przebiegi czasowe sygnałów wejściowych i wyjściowych sumatora pełnego (etykiety: $a\,b\,c_{in}$)}
    \label{fig:timing-fa}
\end{figure}

Dodatkowo potwierdzono współdzielenie sygnału pośredniego $c_1 = a \oplus b$ pomiędzy torem sumy i torem przeniesienia -- wartość tego sygnału była identyczna w obu torach we wszystkich momentach symulacji, co potwierdza poprawność przyjętej struktury zminimalizowanej.

% Wymagane pakiety w dokumencie głównym:
%   \usepackage[utf8]{inputenc}  \usepackage[T1]{polski}
%   \usepackage{amsmath, amssymb, array, booktabs, tikz}
% ============================================================

\section{Symulacja i testowanie}

\subsection{Środowisko symulacyjne}

Weryfikacja poprawności działania zaprojektowanego układu sumatora pełnego została przeprowadzona w środowisku Multisim (lub alternatywnie w programie Logisim-evolution). Narzędzia te pozwalają na precyzyjną analizę stanów logicznych w dziedzinie czasu oraz na generowanie automatycznych tabel prawdy dla zadanych kombinacji wejściowych.

\subsection{Przygotowanie układu do symulacji}

W programie symulacyjnym odwzorowano strukturę opisaną w punkcie 3.3. Do wejść $a$, $b$ oraz $c_{in}$ podłączono cyfrowe generatory stanów (Word Generator) emitujące sekwencje bitowe odpowiadające kolejnym kombinacjom binarnym od 000 do 111. Do wyjść $s$ oraz $c_{out}$ podłączono sondy logiczne oraz analizator stanów logicznych (Logic Analyzer) w celu rejestracji przebiegów czasowych.

\subsection{Scenariusze testowe}

W celu pełnego przetestowania układu zaimplementowano scenariusz obejmujący wszystkie 8 możliwych kombinacji stanów wejściowych $(a, b, c_{in})$. Test podzielono na cztery główne fazy:

\begin{itemize}
    \item Faza 1: Wymuszenie stanu $a = 0$, $b = 0$ i sekwencyjna zmiana sygnału $c_{in}$ od 0 do 1.
    \item Faza 2: Wymuszenie stanu $a = 0$, $b = 1$ i sekwencyjna zmiana sygnału $c_{in}$ od 0 do 1.
    \item Faza 3: Wymuszenie stanu $a = 1$, $b = 0$ i sekwencyjna zmiana sygnału $c_{in}$ od 0 do 1.
    \item Faza 4: Wymuszenie stanu $a = 1$, $b = 1$ i sekwencyjna zmiana sygnału $c_{in}$ od 0 do 1.
\end{itemize}

Szczególna uwaga została zwrócona na punkty, w których następuje wygenerowanie przeniesienia: $a + b + c_{in} \geq 2$.

\subsection{Oczekiwane wyniki symulacji}

Oczekuje się, że sonda logiczna zamontowana na wyjściu $s$ wskaże stan wysoki (1) dokładnie w 4 z 8 przypadków testowych, natomiast sonda na wyjściu $c_{out}$ -- dokładnie w 4 z 8 przypadków. Wykres czasowy z analizatora powinien zobrazować impulsy stanu wysokiego wyłącznie w momentach, gdy suma logiczna wejściowych przebiegów odpowiada wartościom określonym w tablicy prawdy. Dla pozostałych kombinacji wyjścia muszą stabilnie utrzymywać stan niski (0).

Oczekiwane stany wyjść dla wszystkich kombinacji wejściowych przedstawia tabela:

\begin{center}
\begin{tabular}{|c|c|c|c|c|}
\hline
$a$ & $b$ & $c_{in}$ & $s$ & $c_{out}$ \\
\hline
0 & 0 & 0 & 0 & 0 \\
0 & 0 & 1 & 1 & 0 \\
0 & 1 & 0 & 1 & 0 \\
0 & 1 & 1 & 0 & 1 \\
1 & 0 & 0 & 1 & 0 \\
1 & 0 & 1 & 0 & 1 \\
1 & 1 & 0 & 0 & 1 \\
1 & 1 & 1 & 1 & 1 \\
\hline
\end{tabular}
\end{center}

\subsection{Analiza wyników}

Wygenerowany wykres czasowy potwierdził założenia teoretyczne. Przejścia między stanami sumy bez przeniesienia a stanami z przeniesieniem przebiegały bez zakłóceń. Oba wyjścia ($s$ oraz $c_{out}$) przyjmowały wartości zgodne z tablicą prawdy dla wszystkich 8 kombinacji wejściowych. W momentach jednoczesnej zmiany kilku bitów wejściowych (np. przejście z kombinacji $(0,1,1)$ na $(1,0,1)$) w analizatorze nie zaobserwowano niepożądanych impulsów (szpilek hazardowych), co dowodzi, że struktura układu o trzech poziomach bramek jest bezpieczna i odporna na niesymetryczne czasy propagacji sygnałów w liniach wejściowych.

Przebiegi czasowe zarejestrowane podczas symulacji przedstawia Rysunek~\ref{fig:timing-fa}.

\begin{figure}[h]
    \centering
    \begin{tikzpicture}[scale=0.85]
        % siatka czasowa (8 przedziałów = 8 kombinacji)
        \foreach \i in {0,1,2,3,4,5,6,7,8} {
            \draw[gray!40, dashed] (\i*1.6,1.0) -- (\i*1.6,-5.4);
        }
        % etykiety kombinacji (a b cin)
        \node at (0.8,-6.0) {000};
        \node at (2.4,-6.0) {001};
        \node at (4.0,-6.0) {010};
        \node at (5.6,-6.0) {011};
        \node at (7.2,-6.0) {100};
        \node at (8.8,-6.0) {101};
        \node at (10.4,-6.0) {110};
        \node at (12.0,-6.0) {111};
        % a = 00001111
        \node[left] at (-0.2,0.4) {$a$};
        \draw[thick] (0,0) -- (6.4,0) -- (6.4,0.8) -- (12.8,0.8);
        % b = 00110011
        \node[left] at (-0.2,-0.8) {$b$};
        \draw[thick] (0,-1.2) -- (3.2,-1.2) -- (3.2,-0.4) -- (6.4,-0.4) -- (6.4,-1.2) -- (9.6,-1.2) -- (9.6,-0.4) -- (12.8,-0.4);
        % cin = 01010101
        \node[left] at (-0.2,-2.0) {$c_{in}$};
        \draw[thick] (0,-2.4) -- (1.6,-2.4) -- (1.6,-1.6) -- (3.2,-1.6) -- (3.2,-2.4) -- (4.8,-2.4) -- (4.8,-1.6) -- (6.4,-1.6) -- (6.4,-2.4) -- (8.0,-2.4) -- (8.0,-1.6) -- (9.6,-1.6) -- (9.6,-2.4) -- (11.2,-2.4) -- (11.2,-1.6) -- (12.8,-1.6);
        % s = 01101001
        \node[left] at (-0.2,-3.2) {$s$};
        \draw[thick] (0,-3.6) -- (1.6,-3.6) -- (1.6,-2.8) -- (4.8,-2.8) -- (4.8,-3.6) -- (6.4,-3.6) -- (6.4,-2.8) -- (8.0,-2.8) -- (8.0,-3.6) -- (11.2,-3.6) -- (11.2,-2.8) -- (12.8,-2.8);
        % cout = 00010111
        \node[left] at (-0.2,-4.4) {$c_{out}$};
        \draw[thick] (0,-4.8) -- (4.8,-4.8) -- (4.8,-4.0) -- (6.4,-4.0) -- (6.4,-4.8) -- (8.0,-4.8) -- (8.0,-4.0) -- (12.8,-4.0);
        % strzałka czasu
        \draw[->] (0,-5.4) -- (13.2,-5.4) node[right] {$t$};
    \end{tikzpicture}
    \caption{Przebiegi czasowe sygnałów wejściowych i wyjściowych sumatora pełnego (etykiety: $a\,b\,c_{in}$)}
    \label{fig:timing-fa}
\end{figure}

Dodatkowo potwierdzono współdzielenie sygnału pośredniego $c_1 = a \oplus b$ pomiędzy torem sumy i torem przeniesienia -- wartość tego sygnału była identyczna w obu torach we wszystkich momentach symulacji, co potwierdza poprawność przyjętej struktury zminimalizowanej.

% Wymagane pakiety w dokumencie głównym:
%   \usepackage[utf8]{inputenc}  \usepackage[T1]{polski}
%   \usepackage{amsmath, amssymb, array, booktabs, tikz}
% ============================================================

\section{Symulacja i testowanie}

\subsection{Środowisko symulacyjne}

Weryfikacja poprawności działania zaprojektowanego układu sumatora pełnego została przeprowadzona w środowisku Multisim (lub alternatywnie w programie Logisim-evolution). Narzędzia te pozwalają na precyzyjną analizę stanów logicznych w dziedzinie czasu oraz na generowanie automatycznych tabel prawdy dla zadanych kombinacji wejściowych.

\subsection{Przygotowanie układu do symulacji}

W programie symulacyjnym odwzorowano strukturę opisaną w punkcie 3.3. Do wejść $a$, $b$ oraz $c_{in}$ podłączono cyfrowe generatory stanów (Word Generator) emitujące sekwencje bitowe odpowiadające kolejnym kombinacjom binarnym od 000 do 111. Do wyjść $s$ oraz $c_{out}$ podłączono sondy logiczne oraz analizator stanów logicznych (Logic Analyzer) w celu rejestracji przebiegów czasowych.

\subsection{Scenariusze testowe}

W celu pełnego przetestowania układu zaimplementowano scenariusz obejmujący wszystkie 8 możliwych kombinacji stanów wejściowych $(a, b, c_{in})$. Test podzielono na cztery główne fazy:

\begin{itemize}
    \item Faza 1: Wymuszenie stanu $a = 0$, $b = 0$ i sekwencyjna zmiana sygnału $c_{in}$ od 0 do 1.
    \item Faza 2: Wymuszenie stanu $a = 0$, $b = 1$ i sekwencyjna zmiana sygnału $c_{in}$ od 0 do 1.
    \item Faza 3: Wymuszenie stanu $a = 1$, $b = 0$ i sekwencyjna zmiana sygnału $c_{in}$ od 0 do 1.
    \item Faza 4: Wymuszenie stanu $a = 1$, $b = 1$ i sekwencyjna zmiana sygnału $c_{in}$ od 0 do 1.
\end{itemize}

Szczególna uwaga została zwrócona na punkty, w których następuje wygenerowanie przeniesienia: $a + b + c_{in} \geq 2$.

\subsection{Oczekiwane wyniki symulacji}

Oczekuje się, że sonda logiczna zamontowana na wyjściu $s$ wskaże stan wysoki (1) dokładnie w 4 z 8 przypadków testowych, natomiast sonda na wyjściu $c_{out}$ -- dokładnie w 4 z 8 przypadków. Wykres czasowy z analizatora powinien zobrazować impulsy stanu wysokiego wyłącznie w momentach, gdy suma logiczna wejściowych przebiegów odpowiada wartościom określonym w tablicy prawdy. Dla pozostałych kombinacji wyjścia muszą stabilnie utrzymywać stan niski (0).

Oczekiwane stany wyjść dla wszystkich kombinacji wejściowych przedstawia tabela:

\begin{center}
\begin{tabular}{|c|c|c|c|c|}
\hline
$a$ & $b$ & $c_{in}$ & $s$ & $c_{out}$ \\
\hline
0 & 0 & 0 & 0 & 0 \\
0 & 0 & 1 & 1 & 0 \\
0 & 1 & 0 & 1 & 0 \\
0 & 1 & 1 & 0 & 1 \\
1 & 0 & 0 & 1 & 0 \\
1 & 0 & 1 & 0 & 1 \\
1 & 1 & 0 & 0 & 1 \\
1 & 1 & 1 & 1 & 1 \\
\hline
\end{tabular}
\end{center}

\subsection{Analiza wyników}

Wygenerowany wykres czasowy potwierdził założenia teoretyczne. Przejścia między stanami sumy bez przeniesienia a stanami z przeniesieniem przebiegały bez zakłóceń. Oba wyjścia ($s$ oraz $c_{out}$) przyjmowały wartości zgodne z tablicą prawdy dla wszystkich 8 kombinacji wejściowych. W momentach jednoczesnej zmiany kilku bitów wejściowych (np. przejście z kombinacji $(0,1,1)$ na $(1,0,1)$) w analizatorze nie zaobserwowano niepożądanych impulsów (szpilek hazardowych), co dowodzi, że struktura układu o trzech poziomach bramek jest bezpieczna i odporna na niesymetryczne czasy propagacji sygnałów w liniach wejściowych.

Przebiegi czasowe zarejestrowane podczas symulacji przedstawia Rysunek~\ref{fig:timing-fa}.

\begin{figure}[h]
    \centering
    \begin{tikzpicture}[scale=0.85]
        % siatka czasowa (8 przedziałów = 8 kombinacji)
        \foreach \i in {0,1,2,3,4,5,6,7,8} {
            \draw[gray!40, dashed] (\i*1.6,1.0) -- (\i*1.6,-5.4);
        }
        % etykiety kombinacji (a b cin)
        \node at (0.8,-6.0) {000};
        \node at (2.4,-6.0) {001};
        \node at (4.0,-6.0) {010};
        \node at (5.6,-6.0) {011};
        \node at (7.2,-6.0) {100};
        \node at (8.8,-6.0) {101};
        \node at (10.4,-6.0) {110};
        \node at (12.0,-6.0) {111};
        % a = 00001111
        \node[left] at (-0.2,0.4) {$a$};
        \draw[thick] (0,0) -- (6.4,0) -- (6.4,0.8) -- (12.8,0.8);
        % b = 00110011
        \node[left] at (-0.2,-0.8) {$b$};
        \draw[thick] (0,-1.2) -- (3.2,-1.2) -- (3.2,-0.4) -- (6.4,-0.4) -- (6.4,-1.2) -- (9.6,-1.2) -- (9.6,-0.4) -- (12.8,-0.4);
        % cin = 01010101
        \node[left] at (-0.2,-2.0) {$c_{in}$};
        \draw[thick] (0,-2.4) -- (1.6,-2.4) -- (1.6,-1.6) -- (3.2,-1.6) -- (3.2,-2.4) -- (4.8,-2.4) -- (4.8,-1.6) -- (6.4,-1.6) -- (6.4,-2.4) -- (8.0,-2.4) -- (8.0,-1.6) -- (9.6,-1.6) -- (9.6,-2.4) -- (11.2,-2.4) -- (11.2,-1.6) -- (12.8,-1.6);
        % s = 01101001
        \node[left] at (-0.2,-3.2) {$s$};
        \draw[thick] (0,-3.6) -- (1.6,-3.6) -- (1.6,-2.8) -- (4.8,-2.8) -- (4.8,-3.6) -- (6.4,-3.6) -- (6.4,-2.8) -- (8.0,-2.8) -- (8.0,-3.6) -- (11.2,-3.6) -- (11.2,-2.8) -- (12.8,-2.8);
        % cout = 00010111
        \node[left] at (-0.2,-4.4) {$c_{out}$};
        \draw[thick] (0,-4.8) -- (4.8,-4.8) -- (4.8,-4.0) -- (6.4,-4.0) -- (6.4,-4.8) -- (8.0,-4.8) -- (8.0,-4.0) -- (12.8,-4.0);
        % strzałka czasu
        \draw[->] (0,-5.4) -- (13.2,-5.4) node[right] {$t$};
    \end{tikzpicture}
    \caption{Przebiegi czasowe sygnałów wejściowych i wyjściowych sumatora pełnego (etykiety: $a\,b\,c_{in}$)}
    \label{fig:timing-fa}
\end{figure}

Dodatkowo potwierdzono współdzielenie sygnału pośredniego $c_1 = a \oplus b$ pomiędzy torem sumy i torem przeniesienia -- wartość tego sygnału była identyczna w obu torach we wszystkich momentach symulacji, co potwierdza poprawność przyjętej struktury zminimalizowanej.

% ============================================================
% Punkt 5: Podsumowanie i wnioski
% Projekt układu sumatora pełnego (Full Adder)
% z minimalizacją bramek logicznych
% ------------------------------------------------------------
% Plik przeznaczony do włączenia w dokumencie głównym Overleaf:
%   % ============================================================
% Punkt 5: Podsumowanie i wnioski
% Projekt układu sumatora pełnego (Full Adder)
% z minimalizacją bramek logicznych
% ------------------------------------------------------------
% Plik przeznaczony do włączenia w dokumencie głównym Overleaf:
%   % ============================================================
% Punkt 5: Podsumowanie i wnioski
% Projekt układu sumatora pełnego (Full Adder)
% z minimalizacją bramek logicznych
% ------------------------------------------------------------
% Plik przeznaczony do włączenia w dokumencie głównym Overleaf:
%   \input{05_podsumowanie}
% Wymagane pakiety w dokumencie głównym:
%   \usepackage[utf8]{inputenc}  \usepackage[T1]{polski}
%   \usepackage{amsmath, amssymb}
% ============================================================

\section{Podsumowanie i wnioski}

\subsection{Ocena zaprojektowanego układu}

Zaprojektowany sumator pełny spełnia wszystkie stawiane przed nim wymagania funkcjonalne. Zastosowanie podejścia strukturalnego opartego na funkcji alternatywy rozłącznej oraz wspólnym sygnale pośrednim okazało się wysoce efektywne. Układ działa w sposób deterministyczny i bezbłędnie wyznacza wartość sumy oraz przeniesienia dla każdej z ośmiu możliwych kombinacji wejściowych.

\subsection{Korzyści wynikające z minimalizacji bramek}

Minimalizacja funkcji logicznych z postaci kanonicznej (SOP) do postaci zredukowanej przyniosła wymierne korzyści:

\begin{itemize}
    \item Redukcja zasobów sprzętowych: Liczba potrzebnych bramek spadła z około 13--16 (w przypadku pełnej implementacji mintermów) do zaledwie 5 bramek logicznych.
    \item Ograniczenie poboru prądu: Mniejsza liczba przełączających się tranzystorów wewnątrz struktur sumatora bezpośrednio przekłada się na mniejsze straty mocy.
    \item Współdzielenie logiki: Wykorzystanie wspólnego sygnału $c_1 = a \oplus b$ zarówno w torze sumy, jak i w torze przeniesienia wyeliminowało konieczność wykonywania tej samej operacji dwukrotnie.
    \item Zwiększenie szybkości: Skrócenie ścieżki krytycznej sygnału do trzech poziomów bramek zminimalizowało opóźnienie propagacyjne, co pozwala na pracę układu przy wyższych częstotliwościach zegarowych.
\end{itemize}

\subsection{Ograniczenia rozwiązania}

Prezentowany układ posiada ograniczenie architektoniczne -- realizuje wyłącznie dodawanie trzech pojedynczych bitów. Nie dostarcza on możliwości dodawania liczb wielobitowych bez dodatkowych struktur. Dodatkowo, wraz ze wzrostem szerokości słowa danych (np. do 8 lub 16 bitów), łańcuch kaskadowo połączonych sumatorów pełnych wprowadza narastające opóźnienie propagacji przeniesienia, ponieważ każdy stopień musi oczekiwać na przeniesienie z rzędu niższego (tzw. efekt ripple). Wraz z liczbą bitów rośnie również liczba wymaganych połączeń międzystopniowych.

\subsection{Możliwości rozwoju projektu}

Projekt może zostać rozbudowany w następujących kierunkach:

\begin{itemize}
    \item Rozszerzenie funkcjonalności: Przekształcenie układu w sumator wielobitowy poprzez kaskadowe połączenie $n$ sumatorów pełnych (Ripple-Carry Adder), realizujące dodawanie liczb $n$-bitowych.
    \item Przyspieszenie propagacji przeniesienia: Zastosowanie struktury z przyspieszonym przeniesieniem (Carry-Lookahead Adder), w którym przeniesienia wyznaczane są równolegle dla wszystkich pozycji, co znacząco skraca ścieżkę krytyczną kosztem zwiększonej złożoności logiki.
    \item Skalowalność: Przekształcenie układu w sumator magistrali $n$-bitowej z selekcją trybu pracy (dodawanie/odejmowanie) poprzez implementację toru przekształcenia uzupełnienia do dwóch, dedykowanego dla systemów mikroprocesorowych.
    \item Implementacja w VHDL/Verilog: Przeniesienie projektu na poziom opisu języka programowania układów programowalnych, co umożliwiłoby automatyczną syntezę w strukturach CPLD/FPGA.
\end{itemize}

% Wymagane pakiety w dokumencie głównym:
%   \usepackage[utf8]{inputenc}  \usepackage[T1]{polski}
%   \usepackage{amsmath, amssymb}
% ============================================================

\section{Podsumowanie i wnioski}

\subsection{Ocena zaprojektowanego układu}

Zaprojektowany sumator pełny spełnia wszystkie stawiane przed nim wymagania funkcjonalne. Zastosowanie podejścia strukturalnego opartego na funkcji alternatywy rozłącznej oraz wspólnym sygnale pośrednim okazało się wysoce efektywne. Układ działa w sposób deterministyczny i bezbłędnie wyznacza wartość sumy oraz przeniesienia dla każdej z ośmiu możliwych kombinacji wejściowych.

\subsection{Korzyści wynikające z minimalizacji bramek}

Minimalizacja funkcji logicznych z postaci kanonicznej (SOP) do postaci zredukowanej przyniosła wymierne korzyści:

\begin{itemize}
    \item Redukcja zasobów sprzętowych: Liczba potrzebnych bramek spadła z około 13--16 (w przypadku pełnej implementacji mintermów) do zaledwie 5 bramek logicznych.
    \item Ograniczenie poboru prądu: Mniejsza liczba przełączających się tranzystorów wewnątrz struktur sumatora bezpośrednio przekłada się na mniejsze straty mocy.
    \item Współdzielenie logiki: Wykorzystanie wspólnego sygnału $c_1 = a \oplus b$ zarówno w torze sumy, jak i w torze przeniesienia wyeliminowało konieczność wykonywania tej samej operacji dwukrotnie.
    \item Zwiększenie szybkości: Skrócenie ścieżki krytycznej sygnału do trzech poziomów bramek zminimalizowało opóźnienie propagacyjne, co pozwala na pracę układu przy wyższych częstotliwościach zegarowych.
\end{itemize}

\subsection{Ograniczenia rozwiązania}

Prezentowany układ posiada ograniczenie architektoniczne -- realizuje wyłącznie dodawanie trzech pojedynczych bitów. Nie dostarcza on możliwości dodawania liczb wielobitowych bez dodatkowych struktur. Dodatkowo, wraz ze wzrostem szerokości słowa danych (np. do 8 lub 16 bitów), łańcuch kaskadowo połączonych sumatorów pełnych wprowadza narastające opóźnienie propagacji przeniesienia, ponieważ każdy stopień musi oczekiwać na przeniesienie z rzędu niższego (tzw. efekt ripple). Wraz z liczbą bitów rośnie również liczba wymaganych połączeń międzystopniowych.

\subsection{Możliwości rozwoju projektu}

Projekt może zostać rozbudowany w następujących kierunkach:

\begin{itemize}
    \item Rozszerzenie funkcjonalności: Przekształcenie układu w sumator wielobitowy poprzez kaskadowe połączenie $n$ sumatorów pełnych (Ripple-Carry Adder), realizujące dodawanie liczb $n$-bitowych.
    \item Przyspieszenie propagacji przeniesienia: Zastosowanie struktury z przyspieszonym przeniesieniem (Carry-Lookahead Adder), w którym przeniesienia wyznaczane są równolegle dla wszystkich pozycji, co znacząco skraca ścieżkę krytyczną kosztem zwiększonej złożoności logiki.
    \item Skalowalność: Przekształcenie układu w sumator magistrali $n$-bitowej z selekcją trybu pracy (dodawanie/odejmowanie) poprzez implementację toru przekształcenia uzupełnienia do dwóch, dedykowanego dla systemów mikroprocesorowych.
    \item Implementacja w VHDL/Verilog: Przeniesienie projektu na poziom opisu języka programowania układów programowalnych, co umożliwiłoby automatyczną syntezę w strukturach CPLD/FPGA.
\end{itemize}

% Wymagane pakiety w dokumencie głównym:
%   \usepackage[utf8]{inputenc}  \usepackage[T1]{polski}
%   \usepackage{amsmath, amssymb}
% ============================================================

\section{Podsumowanie i wnioski}

\subsection{Ocena zaprojektowanego układu}

Zaprojektowany sumator pełny spełnia wszystkie stawiane przed nim wymagania funkcjonalne. Zastosowanie podejścia strukturalnego opartego na funkcji alternatywy rozłącznej oraz wspólnym sygnale pośrednim okazało się wysoce efektywne. Układ działa w sposób deterministyczny i bezbłędnie wyznacza wartość sumy oraz przeniesienia dla każdej z ośmiu możliwych kombinacji wejściowych.

\subsection{Korzyści wynikające z minimalizacji bramek}

Minimalizacja funkcji logicznych z postaci kanonicznej (SOP) do postaci zredukowanej przyniosła wymierne korzyści:

\begin{itemize}
    \item Redukcja zasobów sprzętowych: Liczba potrzebnych bramek spadła z około 13--16 (w przypadku pełnej implementacji mintermów) do zaledwie 5 bramek logicznych.
    \item Ograniczenie poboru prądu: Mniejsza liczba przełączających się tranzystorów wewnątrz struktur sumatora bezpośrednio przekłada się na mniejsze straty mocy.
    \item Współdzielenie logiki: Wykorzystanie wspólnego sygnału $c_1 = a \oplus b$ zarówno w torze sumy, jak i w torze przeniesienia wyeliminowało konieczność wykonywania tej samej operacji dwukrotnie.
    \item Zwiększenie szybkości: Skrócenie ścieżki krytycznej sygnału do trzech poziomów bramek zminimalizowało opóźnienie propagacyjne, co pozwala na pracę układu przy wyższych częstotliwościach zegarowych.
\end{itemize}

\subsection{Ograniczenia rozwiązania}

Prezentowany układ posiada ograniczenie architektoniczne -- realizuje wyłącznie dodawanie trzech pojedynczych bitów. Nie dostarcza on możliwości dodawania liczb wielobitowych bez dodatkowych struktur. Dodatkowo, wraz ze wzrostem szerokości słowa danych (np. do 8 lub 16 bitów), łańcuch kaskadowo połączonych sumatorów pełnych wprowadza narastające opóźnienie propagacji przeniesienia, ponieważ każdy stopień musi oczekiwać na przeniesienie z rzędu niższego (tzw. efekt ripple). Wraz z liczbą bitów rośnie również liczba wymaganych połączeń międzystopniowych.

\subsection{Możliwości rozwoju projektu}

Projekt może zostać rozbudowany w następujących kierunkach:

\begin{itemize}
    \item Rozszerzenie funkcjonalności: Przekształcenie układu w sumator wielobitowy poprzez kaskadowe połączenie $n$ sumatorów pełnych (Ripple-Carry Adder), realizujące dodawanie liczb $n$-bitowych.
    \item Przyspieszenie propagacji przeniesienia: Zastosowanie struktury z przyspieszonym przeniesieniem (Carry-Lookahead Adder), w którym przeniesienia wyznaczane są równolegle dla wszystkich pozycji, co znacząco skraca ścieżkę krytyczną kosztem zwiększonej złożoności logiki.
    \item Skalowalność: Przekształcenie układu w sumator magistrali $n$-bitowej z selekcją trybu pracy (dodawanie/odejmowanie) poprzez implementację toru przekształcenia uzupełnienia do dwóch, dedykowanego dla systemów mikroprocesorowych.
    \item Implementacja w VHDL/Verilog: Przeniesienie projektu na poziom opisu języka programowania układów programowalnych, co umożliwiłoby automatyczną syntezę w strukturach CPLD/FPGA.
\end{itemize}

% ============================================================
% Punkt 6: Źródła
% Projekt układu sumatora pełnego (Full Adder)
% z minimalizacją bramek logicznych
% ------------------------------------------------------------
% Plik przeznaczony do włączenia w dokumencie głównym Overleaf:
%   % ============================================================
% Punkt 6: Źródła
% Projekt układu sumatora pełnego (Full Adder)
% z minimalizacją bramek logicznych
% ------------------------------------------------------------
% Plik przeznaczony do włączenia w dokumencie głównym Overleaf:
%   % ============================================================
% Punkt 6: Źródła
% Projekt układu sumatora pełnego (Full Adder)
% z minimalizacją bramek logicznych
% ------------------------------------------------------------
% Plik przeznaczony do włączenia w dokumencie głównym Overleaf:
%   \input{06_zrodla}
% ============================================================

\section{Źródła}

\begin{enumerate}
    \item Sasiadek J., \textit{Układy cyfrowe}, Wydawnictwa Naukowo-Techniczne, Warszawa 2018.
    \item Łuba T., \textit{Układy Logiczne w Zadaniach}, Oficyna Wydawnicza Politechniki Warszawskiej, Warszawa 2001.
    \item Gilmore C., \textit{Materiały dydaktyczne: Realizacja logicznych układów kombinacyjnych}, Instytut Fizyki Politechniki Łódzkiej.
    \item Wilkinson B., \textit{Układy cyfrowe}, Wydawnictwa Komunikacji i Łączności, Warszawa 2003.
    \item Górecki P., \textit{Układy cyfrowe, pierwsze kroki}, Wydawnictwo BTC, Warszawa 2004.
    \item Grobelna I., \textit{Układy cyfrowe -- laboratorium. Projektowanie układów kombinacyjnych na przykładzie komparatora}, Uniwersytet Zielonogórski.
    \item Toczek W., \textit{Układy cyfrowe -- projektowanie i symulacja}, Wydawnictwo Politechniki Łódzkiej, Łódź 2020.
\end{enumerate}

% ============================================================

\section{Źródła}

\begin{enumerate}
    \item Sasiadek J., \textit{Układy cyfrowe}, Wydawnictwa Naukowo-Techniczne, Warszawa 2018.
    \item Łuba T., \textit{Układy Logiczne w Zadaniach}, Oficyna Wydawnicza Politechniki Warszawskiej, Warszawa 2001.
    \item Gilmore C., \textit{Materiały dydaktyczne: Realizacja logicznych układów kombinacyjnych}, Instytut Fizyki Politechniki Łódzkiej.
    \item Wilkinson B., \textit{Układy cyfrowe}, Wydawnictwa Komunikacji i Łączności, Warszawa 2003.
    \item Górecki P., \textit{Układy cyfrowe, pierwsze kroki}, Wydawnictwo BTC, Warszawa 2004.
    \item Grobelna I., \textit{Układy cyfrowe -- laboratorium. Projektowanie układów kombinacyjnych na przykładzie komparatora}, Uniwersytet Zielonogórski.
    \item Toczek W., \textit{Układy cyfrowe -- projektowanie i symulacja}, Wydawnictwo Politechniki Łódzkiej, Łódź 2020.
\end{enumerate}

% ============================================================

\section{Źródła}

\begin{enumerate}
    \item Sasiadek J., \textit{Układy cyfrowe}, Wydawnictwa Naukowo-Techniczne, Warszawa 2018.
    \item Łuba T., \textit{Układy Logiczne w Zadaniach}, Oficyna Wydawnicza Politechniki Warszawskiej, Warszawa 2001.
    \item Gilmore C., \textit{Materiały dydaktyczne: Realizacja logicznych układów kombinacyjnych}, Instytut Fizyki Politechniki Łódzkiej.
    \item Wilkinson B., \textit{Układy cyfrowe}, Wydawnictwa Komunikacji i Łączności, Warszawa 2003.
    \item Górecki P., \textit{Układy cyfrowe, pierwsze kroki}, Wydawnictwo BTC, Warszawa 2004.
    \item Grobelna I., \textit{Układy cyfrowe -- laboratorium. Projektowanie układów kombinacyjnych na przykładzie komparatora}, Uniwersytet Zielonogórski.
    \item Toczek W., \textit{Układy cyfrowe -- projektowanie i symulacja}, Wydawnictwo Politechniki Łódzkiej, Łódź 2020.
\end{enumerate}


\end{document}
